\section{Diferenciales en los anillos de característica p}
Los resultados del presente párrafo, de naturaleza más especial y técnica que
los de los §§19, 20 y 22, solo se utilizarán excepcionalmente en el curso del
cap. IV. Su función principal aquí reside en la demostración de tres teoremas del §22
(22.3.3, 22.5.8 y 22.7.3), de los cuales el primero y el último intervienen
esencialmente en la teoría «fina» de los anillos locales noetherianos del
cap. IV, §7.
\oldpage[0\textsubscript{IV-1}]{154}
\subsection{Sistemas de $p$-generadores y $p$-bases}
\label{0.21.1-es}

\begin{env}[21.1.1]
\label{0.21.1.1-es}
Dado un número $p$ que sea o bien $0$ o bien un número primo, diremos que un
anillo $A$ es de característica $p$ si existe un homomorfismo de anillos
$P\to A$, donde $P$ es el cuerpo primo de característica $p$; observemos que
este homomorfismo es entonces único, pues la composición
$\mathbf Z\to P\to A$ es el único homomorfismo de $\mathbf Z$ en $A$.
Si $A\neq0$, equivale a decir que $A$ contiene un cuerpo de característica
$p$, ya que la imagen de $P$ es necesariamente un cuerpo isomorfo a $P$ (y,
por lo demás, el único subcuerpo de $A$ isomorfo a $P$).
\end{env}

\begin{env}[21.1.2]
\label{0.21.1.2-es}
Si $p>0$, decir que $A$ es de característica $p$ equivale a decir que en
$A$ se tiene $p\cdot1=0$, o también $pA=0$. Si $p=0$, decir que $A$ es de
característica $p$ equivale a decir que para todo entero $n\neq0$, $n\cdot1$
es invertible en $A$. Si $A\neq0$, no puede existir más que un único $p$
(primo o $0$) tal que $A$ sea de característica $p$; esto se deduce de la
identidad de Bézout $ap+bq=1$ para dos números primos distintos $p,q$.
En cambio, el anillo reducido a $0$ es de característica $p$ para todo $p$.
\end{env}

\begin{env}[21.1.3]
\label{0.21.1.3-es}
Si $A$ es de característica $p$, lo mismo sucede con toda álgebra sobre $A$.
En particular, para todo ideal primo $\mathfrak p$ de $A$, el cuerpo residual
de $A$ en $\mathfrak p$ es de característica $p$. Recíprocamente, si $p=0$
y, para todo ideal maximal $\mathfrak m$ de $A$, el cuerpo residual de $A$ en
$\mathfrak m$ es de característica $0$, lo mismo sucede con $A$, pues para
todo entero $n\neq0$, $n\cdot1$ es entonces invertible en todos los
$A_{\mathfrak m}$ y, por tanto, también en $A$. En cambio, si $p>0$, un
anillo local puede tener un cuerpo residual de característica $p$ sin ser él
mismo de característica $p$, como muestra el ejemplo del anillo primo
$\mathbf Z_{p\mathbf Z}$ (19.8.3) o del anillo local artiniano
$\mathbf Z/p^n\mathbf Z$ para $n\geq2$, que no contienen cuerpos.

Observemos finalmente que, para un anillo (incluso reducido), los cuerpos
residuales en sus ideales primos pueden tener características distintas, como
muestra el ejemplo de $\mathbf Z$.
\end{env}

\begin{env}[21.1.4]
\label{0.21.1.4-es}
En todo lo que sigue suponemos fijado un número primo $p$ y todos los anillos
se supondrán de característica $p$, salvo mención expresa en contrario. Para
un anillo $A$ de este tipo, la aplicación $x\mapsto x^p$ es un endomorfismo de
$A$, que denotamos por $F_A$; si $A$ es reducido, $F_A$ es inyectivo. Ponemos
$A^p=F_A(A)$ (subanillo de $A$ formado por los $x^p$, $x\in A$); podemos
considerar naturalmente $A$ como una $A^p$-álgebra.

También podemos considerar $A$ como una $A$-álgebra mediante el homomorfismo
$F_A:x\mapsto x^p$ de $A$ en $A$; en otras palabras, se trata del álgebra
$A^{(p)}$ para la cual el producto $\lambda\cdot x$ de $x\in A$ por un
escalar $\lambda\in A$ es el producto $\lambda^p x$ en el anillo $A$. Es
claro que, para todo homomorfismo de anillos $u:A\to B$, el par $(u,u)$ es un
dihomomorfismo de álgebras $A^{(p)}\to B^{(p)}$.
Para todo $A$-módulo $E$, ponemos $E^{(p)}=E\otimes_A A^{(p)}$, donde
$A^{(p)}$ se considera como un $(A,A)$-bimódulo, siendo la estructura de
$A$-módulo por la izquierda la que acabamos de definir y la estructura de
$A$-módulo por la derecha la definida por la multiplicación en $A$; así,
para $x\in E$, $\alpha,\beta\in A$, se tiene
\[
  \alpha(x\otimes\beta)=x\otimes(\alpha\beta),\qquad
  (\alpha x)\otimes\beta=x\otimes\alpha^p\beta.
\]
\end{env}

\oldpage[0\textsubscript{IV-1}]{155}
Poniendo $x^{(p)}=x\otimes1$, se tiene, por tanto,
$(\alpha x)^{(p)}=\alpha^p x^{(p)}$. Cuando $F_A$ es inyectivo, podemos
identificar $A^{(p)}$ con el anillo $A$ considerado como álgebra sobre su
subanillo $A^p$.

\begin{proposition}[21.1.5]
\label{0.21.1.5-es}
Sean $A$, $B$ dos anillos y $u:A\to B$ un homomorfismo.
\begin{enumerate}
  \item[(i)] Toda $A$-derivación de $B$ en un $B$-módulo $L$ se anula en
  $A[B^p]$ (por tanto es una $A[B^p]$-derivación).
  \item[(ii)] Para toda sub-$A$-álgebra $A'$ de $A[B^p]$, si
  $j:A'\to B$ es la inyección canónica, el homomorfismo canónico
  \[
  j_{B/A'/A}:\Omega^1_{B/A}\longrightarrow\Omega^1_{B/A'}
  \]
  es biyectivo.
  \item[(iii)] Supongamos que existe un entero $s\geq0$ tal que
  $u(A)\supset B^{p^s}$. Entonces, en el anillo $B\otimes_A B$,
  $\mathfrak J_{B/A}$ es un ideal nil.
\end{enumerate}
\end{proposition}

\begin{proof}
(i) Por inducción a partir de (20.1.1.1), se deduce, para toda
$A$-derivación $D:B\to L$, que $D(x^n)=nx^{n-1}D(x)$, y en particular
$D(x^p)=0$.

(ii) En vista de (20.4.8), la afirmación (ii) no es más que una traducción
de (i), que se escribe
$\operatorname{Der}_{A[B^p]}(B,L)=\operatorname{Der}_A(B,L)$ para todo
$B$-módulo $L$.

(iii) Para todo $x\in B$, se tiene
$(x\otimes1-1\otimes x)^{p^s}=x^{p^s}\otimes1-1\otimes x^{p^s}=0$,
puesto que $x^{p^s}\in u(A)$. La conclusión resulta de que los elementos
$1\otimes x-x\otimes1$ generan $\mathfrak J_{B/A}$ (20.4.4).
\end{proof}

Se deduce inmediatamente de (21.1.5) que, para todo par de homomorfismos de
anillos $A\to B\to C$, se tiene
\begin{equation}
\tag{21.1.5.1}
\label{0.21.1.5.1-es}
\Upsilon_{C/B/A}=\Upsilon_{C/B/A'}
\end{equation}
para toda sub-$A$-álgebra $A'\subset A[B^p]$ de $B$.

Por otra parte, (21.1.5) muestra en particular que, para los módulos de
diferenciales «absolutas»,
\begin{equation}
\tag{21.1.5.2}
\label{0.21.1.5.2-es}
\Omega^1_A=\Omega^1_{A/A^p}.
\end{equation}

\begin{corollary}[21.1.6]
\label{0.21.1.6-es}
Supongamos que $B$ es una $A$-álgebra de tipo finito y que existe un entero
$s\geq0$ tal que $u(A)\supset B^{p^s}$. Si $\Omega^1_{B/A}=0$, $u$ es
sobreyectivo.
\end{corollary}

\begin{proof}
\oldpage[0\textsubscript{IV-1}]{156}
Según (21.1.5, (iii)), $\mathfrak J_{B/A}$ es un ideal nil; además, en
virtud de (20.4.4) y de la hipótesis, $\mathfrak J_{B/A}$ es un ideal de tipo
finito y, por tanto, es nilpotente; como la relación $\Omega^1_{B/A}=0$
significa que $\mathfrak J_{B/A}=\mathfrak J_{B/A}^2$, se concluye que
$\mathfrak J_{B/A}=0$, o, equivalentemente, que
$\pi:B\otimes_A B\to B$ es biyectivo.
Por otra parte, todo elemento de $B$ cuya potencia $p^s$-ésima pertenece a
$u(A)$ es entero; por tanto, $B$ es entero sobre $A$ y, al ser de tipo finito,
es una $A$-álgebra finita. Se reduce así la cuestión a demostrar el lema
siguiente (en el que no
se supone que los anillos considerados sean de característica $p$):
\end{proof}

\begin{lemma}[21.1.6.1]
\label{0.21.1.6.1-es}
Sean $R$ un anillo y $S$ una $R$-álgebra finita; si el homomorfismo canónico
$\pi:S\otimes_R S\to S$ es biyectivo, entonces el homomorfismo estructural
$u:R\to S$ es sobreyectivo.
\end{lemma}

\begin{proof}
Basta demostrar que, para todo ideal maximal $\mathfrak m$ de $R$, si se
pone $T=R-\mathfrak m$, el homomorfismo
$u_{\mathfrak m}:R_{\mathfrak m}\to T^{-1}S$ es sobreyectivo (Bourbaki,
\emph{Alg. comm.}, cap.~II, §~3, n$^\circ$~3, teorema~1); ahora bien, la
hipótesis implica que el homomorfismo
$(T^{-1}S)\otimes_{R_{\mathfrak m}}(S/\mathfrak mS)\to T^{-1}S$
es biyectivo (0$_I$, 1.3.4), y como $T^{-1}S$ es una
$R_{\mathfrak m}$-álgebra finita, se ve que podemos limitarnos al caso en
que $R$ es un anillo local.

Denotando todavía por $\mathfrak m$ su ideal maximal, basta demostrar que
$u\otimes1:R/\mathfrak m\to S/\mathfrak mS$ es sobreyectivo, en virtud del
lema de Nakayama ($S$ es un $R$-módulo de tipo finito); como el homomorfismo
canónico
$(S/\mathfrak mS)\otimes_{R/\mathfrak m}(S/\mathfrak mS)\to S/\mathfrak mS$
es biyectivo y $S/\mathfrak mS$ es una $(R/\mathfrak m)$-álgebra finita,
finalmente nos reducimos al caso en que $R$ es un cuerpo; pero entonces los
rangos de $S\otimes_R S$ y de $S$ sobre $R$ solo pueden ser iguales si $S$
es de rango $0$ o $1$, es decir, si $u:R\to S$ es sobreyectivo.
\end{proof}

\begin{proposition}[21.1.7]
\label{0.21.1.7-es}
Sean $A$ un anillo, $B$ una $A$-álgebra, $(x_\alpha)_{\alpha\in I}$ una
familia de elementos de $B$. Consideremos las propiedades siguientes:
\begin{enumerate}
  \item[a)] $B=A[B^p,(x_\alpha)]$, es decir, la $A[B^p]$-álgebra $B$ está
  generada por la familia $(x_\alpha)_{\alpha\in I}$.
  \item[b)] El $A[B^p]$-módulo $B$ está generado por los monomios
  $\prod_\alpha x_\alpha^{n(\alpha)}$, donde $(n(\alpha))_{\alpha\in I}$ es
  una familia de enteros de soporte finito tal que $0\leq n(\alpha)<p$ para
  todo $\alpha\in I$.
  \item[c)] El $B$-módulo $\Omega^1_{B/A}$ está generado por los
  $d_{B/A}(x_\alpha)$ $(\alpha\in I)$.
\end{enumerate}
Entonces las propiedades a) y b) son equivalentes y entrañan c); además, si
$B$ es una $A[B^p]$-álgebra de tipo finito, c) es equivalente a a) y b).
\end{proposition}

\begin{proof}
Está claro que b) implica a), y, recíprocamente, a) implica b), pues todo
monomio $\prod_\alpha x_\alpha^{m(\alpha)}$, donde los $m(\alpha)$ son
enteros $\geq0$ (formando una familia de soporte finito), puede escribirse
$(\prod_\alpha x_\alpha^{q(\alpha)})^p\prod_\alpha x_\alpha^{r(\alpha)}$,
dividiendo cada $m(\alpha)$ por $p$, lo que da
$m(\alpha)=p q(\alpha)+r(\alpha)$ con $0\leq r(\alpha)<p$.
El hecho de que a) entrañe c) se deduce de (21.1.5, (ii)) y de (20.4.7).

Supongamos recíprocamente que c) se verifica y que $B$ es una
$A[B^p]$-álgebra de tipo finito; sea $B'$ la sub-$A[B^p]$-álgebra de $B$
generada por los $x_\alpha$; en la sucesión exacta (20.5.7.1)
\[
\Omega^1_{B'/A[B^p]}\otimes_{B'}B\longrightarrow
\Omega^1_{B/A[B^p]}\longrightarrow\Omega^1_{B/B'}\longrightarrow0
\]
la hipótesis entraña que la flecha de la izquierda es sobreyectiva (teniendo
en cuenta (21.1.5, (ii))); por tanto, $\Omega^1_{B/B'}=0$, y como
$B^p\subset B'\subset B$ y $B$ es una $B'$-álgebra de tipo finito, se tiene
necesariamente $B'=B$ por (21.1.6).
\end{proof}

\begin{remark}[21.1.8]
\label{0.21.1.8-es}
Cuando $B$ es un cuerpo, demostraremos la equivalencia de las propiedades a),
b) y c) sin hipótesis de finitud (21.4.5).
\end{remark}

\begin{definition}[21.1.9]
\label{0.21.1.9-es}
Sean $A$ un anillo (de característica $p$), $B$ una $A$-álgebra. Diremos
que una familia $(x_\alpha)_{\alpha\in I}$ de elementos de $B$ es
$p$-libre sobre $A$ (resp. un sistema de $p$-generadores de $B$ sobre $A$,
resp. una $p$-base de $B$ sobre $A$) si la familia de monomios
$\prod_\alpha x_\alpha^{n(\alpha)}$ ($0\leq n(\alpha)<p$,
$(n(\alpha))_{\alpha\in I}$ de soporte finito) es una familia libre (resp.
un sistema de generadores, resp. una base) en el $A[B^p]$-módulo $B$.
\end{definition}

Se tienen definiciones correspondientes para un conjunto $M$ de elementos de
$B$, considerando la familia definida por la inyección canónica $M\to B$.
Cuando se toma $A$ igual al cuerpo primo $\mathbf F_p$ (en cuyo caso
$A[B^p]=B^p$), se omite la mención de $A$ en la definición anterior (o se
dice también que una familia es «absolutamente» $p$-libre, resp. un sistema
«absoluto» de $p$-generadores, resp. una $p$-base «absoluta»).

\oldpage[0\textsubscript{IV-1}]{157}
\begin{lemma}[21.1.10]
\label{0.21.1.10-es}
Sean $C$ una sub-$A$-álgebra de $B$ tal que $B^p\subset C$, $M$ una
parte de $C$ y $N$ una parte de $B$.
\begin{enumerate}
  \item[(i)] Si $M$ es un sistema de $p$-generadores de $C$ sobre $A$ y $N$
  es un sistema de $p$-generadores de $B$ sobre $C$, entonces $M\cup N$ es un
  sistema de $p$-generadores de $B$ sobre $A$.
  \item[(ii)] Supongamos que $M$ es una $p$-base de $C$ sobre $A$; entonces,
  para que $N$ sea $p$-libre sobre $C$, es necesario y suficiente que $M\cup N$
  sea $p$-libre sobre $A$.
\end{enumerate}
\end{lemma}

\begin{proof}
Podemos limitarnos al caso en que $B^p\subset A\subset C\subset B$. Entonces
(i) es un caso particular del hecho de que, si $P$ (resp. $Q$) es un sistema
de generadores del $A$-módulo $C$ (resp. del $C$-módulo $B$), el conjunto de
los productos $xy$, donde $x\in P$ y $y\in Q$, es un sistema de generadores
del $A$-módulo $B$.

Conservando las mismas notaciones, si $P$ es una base del $A$-módulo $C$,
decir que $Q$ es una familia libre sobre $C$ significa que la relación
$\sum_{\lambda,\mu}a_{\lambda\mu}x_\lambda y_\mu=0$, donde
$a_{\lambda\mu}\in A$, $x_\lambda\in P$, $y_\mu\in Q$ (siendo los
$x_\lambda$ (resp. los $y_\mu$) distintos dos a dos), equivale a
$a_{\lambda\mu}=0$ para todo par $(\lambda,\mu)$, o también a
$\sum_\lambda a_{\lambda\mu}x_\lambda=0$ para todo $\mu$; de ahí la
afirmación (ii).
\end{proof}

\subsection{$p$-bases y lisitud formal}
\label{0.21.2-es}

\begin{theorem}[21.2.1]
\label{0.21.2.1-es}
Sean $B$ un anillo, $A$ un subanillo de $B$ tal que $B^p\subset A\subset B$,
$(x_\alpha)_{\alpha\in I}$ una $p$-base de $B$ sobre $A$.
Sean $E$ una $A$-álgebra, $u:A\to E$ el homomorfismo estructural,
$q:E\to B$ un $A$-homomorfismo sobreyectivo, $\mathfrak R$ su núcleo, y
supongamos que $\mathfrak R^p=0$. Entonces:
\begin{enumerate}
  \item[(i)] Para que exista un $A$-homomorfismo $v:B\to E$ inverso por la
  derecha del homomorfismo $q:E\to B$, es necesario y suficiente que
  $u(q(z)^p)=z^p$ para todo $z\in E$.
  \item[(ii)] Cuando se satisface la condición de (i), para toda familia
  $(z_\alpha)_{\alpha\in I}$ de elementos de $E$ tal que
  $q(z_\alpha)=x_\alpha$ para todo $\alpha\in I$, existe un único
  $A$-homomorfismo $v:B\to E$ tal que $v(x_\alpha)=z_\alpha$ para todo
  $\alpha\in I$, y $v$ es inverso por la derecha de $q$.
\end{enumerate}
\end{theorem}

\begin{proof}
Si existe un $A$-homomorfismo $v:B\to E$, se debe tener $v(a)=u(a)$ para
todo $a\in A\subset B$, y por tanto $v(q(z)^p)=u(q(z)^p)$ para todo $z\in E$,
ya que $B^p\subset A$. Pero se tiene $v(q(z)^p)=(v(q(z)))^p$ y, por definición,
$v(q(z))\equiv z\pmod{\mathfrak R}$ si $q\circ v=1_B$, de modo que
$(v(q(z)))^p=z^p$ puesto que $\mathfrak R^p=0$; de ahí la necesidad de (i).
La suficiencia de la condición (i) resultará de (ii).

Ahora bien, bajo las hipótesis de (ii), la unicidad de $v$ es evidente, pues
los monomios $\prod_\alpha x_\alpha^{n(\alpha)}$ generan el $A$-módulo $B$;
como estos monomios forman además una \emph{base} del $A$-módulo $B$, existe
una aplicación $A$-\emph{lineal} $v$ de $B$ en $E$ tal que
$v(\prod_\alpha x_\alpha^{n(\alpha)})=\prod_\alpha z_\alpha^{n(\alpha)}$
para toda familia $(n(\alpha))_{\alpha\in I}$ de soporte finito con
$0\leq n(\alpha)<p$ para todo $\alpha$. Queda por ver que $v$ es un
homomorfismo de anillos. Ahora bien, se puede escribir
$(\prod_\alpha x_\alpha^{m(\alpha)})(\prod_\alpha x_\alpha^{n(\alpha)})
=a\prod_\alpha x_\alpha^{r(\alpha)}$, donde
$r(\alpha)=m(\alpha)+n(\alpha)$ si $m(\alpha)+n(\alpha)<p$,
$r(\alpha)=m(\alpha)+n(\alpha)-p$ en caso contrario, y $a\in A$
\oldpage[0\textsubscript{IV-1}]{158}
es el producto de los $x_\alpha^p$ para los $\alpha\in I$ tales que
$m(\alpha)+n(\alpha)\geq p$; se trata de comprobar que
$u(a)=\prod_\alpha z_\alpha^p$; pero como $x_\alpha=q(z_\alpha)$, esto se
deduce de la hipótesis.
\end{proof}

\begin{corollary}[21.2.2]
\label{0.21.2.2-es}
Sean $B$ un anillo, $A$ un subanillo de $B$ tal que $B^p\subset A\subset B$.
Si existe una $p$-base de $B$ sobre $A$, $B$ es una $A$-álgebra formalmente
lisa respecto de $B^p$ (para las topologías discretas).
\end{corollary}

\begin{proof}
Sean $E$ una $A$-extensión de $B$ por un $B$-módulo $L$, $q:E\to B$ la
aumentación, $u:A\to E$ el homomorfismo estructural. Decir que $E$ es
$B^p$-trivial significa que existe un homomorfismo de anillos $v:B\to E$
tal que $q(v(b))=b$ y $v(b^p)=u(b^p)$ para todo $b\in B$. Se deduce, para
$z\in E$, que $u(q(z)^p)=v(q(z)^p)=(v(q(z)))^p$; pero, en virtud de la
relación $L^2=0$, también se tiene $L^p=0$, y como $v(q(z))-z\in L$, se
tiene $(v(q(z)))^p=z^p$. Por tanto se satisface la condición de (21.2.1),
(i), y $E$ es $B^p$-trivial.
\end{proof}

\begin{corollary}[21.2.3]
\label{0.21.2.3-es}
Sean $B$ un anillo, $A$ un subanillo de $B$ tal que $B^p\subset A\subset B$,
y $\left(x_\alpha\right)_{\alpha\in I}$ una $p$-base de $B$ sobre $A$. Sea
$L$ un $B$-módulo. Entonces:
\begin{enumerate}
  \item[(i)] Para que una derivación $D$ de $A$ en $L$ se prolongue a una
  derivación de $B$ en $L$, es necesario y suficiente que $D$ se anule en
  $B^p$.
  \item[(ii)] Si $D$ se anula en $B^p$, entonces, para toda familia
  $(y_\alpha)_{\alpha\in I}$ de elementos de $L$, existe una única derivación
  $D'$ de $B$ en $L$ que prolonga $D$ y tal que
  $D'(x_\alpha)=y_\alpha$ para todo $\alpha\in I$.
\end{enumerate}
\end{corollary}

\begin{proof}
Dada una derivación $D$ de $A$ en $L$, consideremos el anillo $E=D_B(L)$ y
el homomorfismo $u:A\to E$ definido por $u(a)=(a,D(a))$; un
$A$-homomorfismo $v:B\to E$ inverso por la derecha del homomorfismo canónico
$q:E\to B$ es entonces de la forma $x\mapsto(x,D'(x))$, donde $D'$ es una
derivación de $B$ en $L$ que prolonga $D$ (20.1.5); como
$u(q(z)^p)=(q(z)^p,D(q(z)^p))$ para $z\in E$, el corolario se deduce
inmediatamente de (21.2.1) aplicado a $E$.
\end{proof}

Se deduce de (21.2.3) que la sucesión
\begin{equation}
\tag{21.2.3.1}
\label{0.21.2.3.1-es}
0\to\operatorname{Der}_A(B,L)\to\operatorname{Der}_{B^p}(B,L)
\to\operatorname{Der}_{B^p}(A,L)\to0
\end{equation}
(cf. (20.2.2.1)) es exacta, y que la aplicación
\begin{equation}
\tag{21.2.3.2}
\label{0.21.2.3.2-es}
D'\longmapsto(D'(x_\alpha))_{\alpha\in I}
\end{equation}
es un isomorfismo del $B$-módulo $\operatorname{Der}_A(B,L)$ sobre el
$B$-módulo producto $L^I$ (tomando $D=0$ en (21.2.3)).

\begin{corollary}[21.2.4]
\label{0.21.2.4-es}
Bajo las hipótesis de (21.2.3), la sucesión
\begin{equation}
\tag{21.2.4.1}
\label{0.21.2.4.1-es}
0\to\Omega^1_{A/B^p}\otimes_A B\to\Omega^1_{B/B^p}
\to\Omega^1_{B/A}\to0
\end{equation}
es exacta y escindida, y la familia $(d_{B/A}(x_\alpha))_{\alpha\in I}$ es
una base del $B$-módulo $\Omega^1_{B/A}$.
\end{corollary}

\begin{proof}
Esto se deduce inmediatamente de (21.2.3) y de la fórmula
$\operatorname{Der}_A(B,L)=\operatorname{Hom}_B(\Omega^1_{B/A},L)$ (20.4.8).
\end{proof}

\begin{corollary}[21.2.5]
\label{0.21.2.5-es}
Sean $A$ un anillo, $B$ una $A$-álgebra y $(x_\alpha)_{\alpha\in I}$ una
$p$-base de $B$ sobre $A$. Entonces $(d_{B/A}(x_\alpha))_{\alpha\in I}$ es
una base del $B$-módulo $\Omega^1_{B/A}$.
\end{corollary}

\begin{proof}
Usando (21.1.5), (ii), nos reducimos al caso en que $A=A[B^p]$ y basta
entonces aplicar (21.2.4).
\end{proof}

\oldpage[0\textsubscript{IV-1}]{159}
\begin{env}[21.2.6]
\label{0.21.2.6-es}
Sean $A$, $B$ dos anillos, $u:A\to B$ un homomorfismo de anillos; se tiene
(con las notaciones de (21.1.4)) el diagrama conmutativo
\begin{equation}
\tag{21.2.6.1}
\label{0.21.2.6.1-es}
\xymatrix{
 A^{(p)} \ar[r]^{u} & B^{(p)}\\
 A \ar[u]^{F_A} \ar[r]^{u} & B \ar[u]_{F_B}
}
\end{equation}
De él se deduce canónicamente un homomorfismo de $B$-álgebras
\begin{equation}
\tag{21.2.6.2}
\label{0.21.2.6.2-es}
u\otimes1:A^{(p)}\otimes_A B\longrightarrow B^{(p)}
\end{equation}
cuya imagen es el anillo $A[B^p]$ (que es una $B$-álgebra para el
homomorfismo $F_B:B\to B^p$).
\end{env}

\begin{theorem}[21.2.7]
\label{0.21.2.7-es}
Sean $A$ un anillo, $B$ una $A$-álgebra, $u:A\to B$ el homomorfismo
estructural. Supongamos verificadas las condiciones siguientes:
\begin{enumerate}
  \item[(i)] El homomorfismo (21.2.6.2) deducido de $u$ es inyectivo.
  \item[(ii)] $B$ admite una $p$-base $(x_\alpha)_{\alpha\in I}$ respecto de
  $A$.
\end{enumerate}
Entonces $B$ es una $A$-álgebra formalmente lisa (para las topologías
discretas). Más precisamente, dotemos $A$ y $B$ de las topologías discretas;
sean $E$ una $A$-álgebra topológica admisible (0$_I$, 7.1.2),
$\mathfrak R$ un ideal de definición de $E$, $C=E/\mathfrak R$,
$v:B\to C$ un $A$-homomorfismo, $q:E\to C$ la aumentación. Entonces, para
toda familia $(z_\alpha)_{\alpha\in I}$ de elementos de $E$ tal que
$q(z_\alpha)=v(x_\alpha)$ para todo $\alpha\in I$, existe un único
$A$-homomorfismo $w:B\to E$ tal que $v=q\circ w$ y
$w(x_\alpha)=z_\alpha$ para todo $\alpha\in I$.
\end{theorem}

\begin{proof}
Consideremos primero el caso en que $E$ es discreto, de modo que $\mathfrak R$
es nilpotente. Podemos limitarnos al caso en que $v$ es sobreyectivo; además,
el razonamiento de (19.4.3) permite suponer $\mathfrak R^2=0$ y, por tanto,
$\mathfrak R^p=0$. Finalmente, considerando la imagen inversa por $v$ de la
extensión $E$ de $C$ por $\mathfrak R$, podemos limitarnos al caso en que $C=B$
y $v$ es la identidad (19.4.4). Como el homomorfismo de anillos
$F_E:z\rightsquigarrow z^p$ se anula en $\mathfrak R$, se factoriza como
$E\xrightarrow q B\xrightarrow{F'}E$, y $F'$, considerado como homomorfismo de
$B$ en $E^{(p)}$, es un \emph{$A$-homomorfismo} por definición de la estructura de
$A$-álgebra de $E^{(p)}$ mediante el homomorfismo compuesto
$A\xrightarrow r E\xrightarrow{F_E}E^{(p)}$, donde $r:A\to E$ es el homomorfismo
estructural de la $A$-álgebra $E$; por otra parte,
$r:A^{(p)}\to E^{(p)}$ es también un $A$-homomorfismo. Existe, pues, un único
$A$-homomorfismo $f:A^{(p)}\otimes_A B\to E^{(p)}$ tal que las composiciones de
$f$ con los homomorfismos canónicos $A^{(p)}\to A^{(p)}\otimes_A B$ y
$B\to A^{(p)}\otimes_A B$ sean, respectivamente, $r$ y $F'$. Ahora bien, por
hipótesis, podemos identificar $A^{(p)}\otimes_A B$ con $A[B^p]$, identificándose
los homomorfismos canónicos $A^{(p)}\to A^{(p)}\otimes_A B$ y
$B\to A^{(p)}\otimes_A B$ respectivamente con $u$ y con $F_B$. Podemos, pues,
considerar ahora $E$ como una $A[B^p]$-álgebra mediante el homomorfismo $f$, y,
por construcción, tenemos $f(q(z)^p)=z^p$ para todo $z\in E$; se cumplen, por
tanto, las condiciones de aplicación de (21.2.1), de donde se deduce el
teorema en este caso.

Para pasar al caso general, consideremos un sistema fundamental
$(\mathfrak J_\lambda)$ de ideales abiertos de $E$, y pongamos
$E_\lambda=E/\mathfrak J_\lambda$,
$\mathfrak R_\lambda=(\mathfrak R+\mathfrak J_\lambda)/\mathfrak J_\lambda$,
$C_\lambda=E_\lambda/\mathfrak R_\lambda=E/(\mathfrak R+\mathfrak J_\lambda)$;
como $E$ es admisible, tenemos $E=\varprojlim E_\lambda$ (0$_I$, 7.2.2).
Para cada par $(\alpha,\lambda)$, denotemos por $z_{\alpha\lambda}$ la imagen
canónica de $z_\alpha$ en $E_\lambda$, por $p_\lambda:C\to C_\lambda$ el
homomorfismo canónico,
\oldpage[0\textsubscript{IV-1}]{160}
y por $q_\lambda$ el homomorfismo canónico $E_\lambda\to C_\lambda$.
Como, por hipótesis, $\mathfrak R_\lambda$ es nilpotente en $E_\lambda$, la primera
parte de la demostración prueba la existencia y unicidad de un
$A$-homomorfismo $w_\lambda:B\to E_\lambda$ tal que
$p_\lambda\circ v=q_\lambda\circ w_\lambda$ y
$w_\lambda(x_\alpha)=z_{\alpha\lambda}$ para todo $\alpha$.
La unicidad de los $w_\lambda$ muestra entonces inmediatamente que $(w_\lambda)$
es un sistema proyectivo de homomorfismos, y $w=\varprojlim w_\lambda$ responde
a la cuestión.
\end{proof}

\begin{remark}[21.2.8]
\label{0.21.2.8-es}
La verificación de la existencia y unicidad del homomorfismo $w:B\to E$ tal
que $w(x_\alpha)=z_\alpha$ para todo $\alpha$ se deduce directamente del
hecho de que $B$ es una $A$-álgebra formalmente lisa y de que los
$d_{B/A}(x_\alpha)$ forman una base del $B$-módulo $\Omega^1_{B/A}$ (21.2.4),
sin hacer intervenir el hecho de que se trata de una $p$-base (de modo que el
resultado es válido sin suponer que los anillos sean de característica $p$).
En efecto, podemos limitarnos al caso en que $\mathfrak R^2=0$; como
$\operatorname{Der}_A(B,\mathfrak R)=\operatorname{Hom}_B(\Omega^1_{B/A},\mathfrak R)$
(20.4.8), existe una única $A$-derivación $D$ de $B$ en $\mathfrak R$ tal que
$D(x_\alpha)=z'_\alpha$ para toda familia $(z'_\alpha)$ de elementos de
$\mathfrak R$; la conclusión se deduce de (20.1.1).
\end{remark}

\subsection{$p$-bases y módulos de imperfección}
\label{0.21.3-es}

\begin{env}[21.3.1]
\label{0.21.3.1-es}
Sean $A$, $B$ dos anillos (de característica $p$), $i:A\to B$,
$j:B\to A$ dos homomorfismos de anillos tales que
\begin{align}
\tag{21.3.1.1}\label{0.21.3.1.1-es}
j(i(a))&=a^p\quad\text{para todo }a\in A,\\
\tag{21.3.1.2}\label{0.21.3.1.2-es}
i(j(b))&=b^p\quad\text{para todo }b\in B.
\end{align}
Con frecuencia $i$ será inyectivo, de modo que $A$ se identificará mediante
$i:A\to B$ con un subanillo de $B$; hecha esta identificación, la existencia
de $j$ implica que $B^p\subset A$, y $j$ se identifica entonces con $F_B$.
\end{env}

\begin{env}[21.3.2]
\label{0.21.3.2-es}
Si $h:A^p\to A$ es la inyección canónica, se tiene, según (21.3.1.1) y
(21.3.1.2), un diagrama conmutativo
\begin{equation}
\tag{21.3.2.1}\label{0.21.3.2.1-es}
\xymatrix{
 B \ar[r]^{j} & A\\
 A \ar[u]^{i} \ar[r]^{F_A} & A^p \ar[u]_{h}
}
\end{equation}
por lo que el par $(j,F_A)$ puede considerarse un dihomomorfismo del
$A$-álgebra $B$ (para $i$) en el $A^p$-álgebra $A$ (para $h$).
De ello se deduce un homomorfismo canónico de $A$-módulos
\begin{equation}
\tag{21.3.2.2}\label{0.21.3.2.2-es}
\pi_{B/A}:\Omega^1_{B/A}\otimes_B A_{[j]}\longrightarrow
\Omega^1_{A/A^p}=\Omega^1_A
\end{equation}
(cf. (20.5.4); la identificación de $\Omega^1_{A/A^p}$ con el módulo de
diferenciales absolutas $\Omega^1_A$ proviene de (21.1.5), (ii)).
\begin{align}
\tag{21.3.2.3}\label{0.21.3.2.3-es}
\Theta_{B/A}&=\Omega^1_{B/A}\otimes_B A_{[j]},\\
\tag{21.3.2.4}\label{0.21.3.2.4-es}
\Xi_{B/A}&=\operatorname{Ker}(\pi_{B/A}).
\end{align}
\oldpage[0\textsubscript{IV-1}]{161}
Por tanto se tiene la sucesión exacta
\begin{equation}
\tag{21.3.2.5}\label{0.21.3.2.5-es}
0\to\Xi_{B/A}\to\Theta_{B/A}\xrightarrow{\pi_{B/A}}\Omega^1_A
\end{equation}
(se observará que estas notaciones pueden inducir a confusión, pues
$\Theta_{B/A}$ y $\Xi_{B/A}$ dependen no solo de $A$ y $B$, sino también de
$i$ y $j$).
\end{env}

\begin{env}[21.3.3]
\label{0.21.3.3-es}
Como, en virtud de (21.3.1.2), se tiene $F_B=i\circ j$, se puede escribir,
para todo $B$-módulo $M$ (cf. (21.1.4)),
\begin{equation}
\tag{21.3.3.1}\label{0.21.3.3.1-es}
M^{(p)}=M\otimes_B B^{(p)}=(M\otimes_B A_{[j]})\otimes_A B_{[i]},
\end{equation}
de donde, en particular,
\begin{equation}
\tag{21.3.3.2}\label{0.21.3.3.2-es}
(\Omega^1_{B/A})^{(p)}=\Theta_{B/A}\otimes_A B_{[i]},
\end{equation}
y, por tanto, de (21.3.2.2) se deduce un homomorfismo canónico de
$B$-módulos
\begin{equation}
\tag{21.3.3.3}\label{0.21.3.3.3-es}
\pi_{B/A}\otimes1_B:(\Omega^1_{B/A})^{(p)}\longrightarrow
\Omega^1_A\otimes_A B_{[i]}.
\end{equation}
Teniendo en cuenta (20.5.4.2), la imagen de este homomorfismo es el $B$-módulo
generado por los elementos $d_A(j(b))\otimes1$ para $b\in B$. Su imagen por
el homomorfismo canónico $i_{B/A}:\Omega^1_A\otimes_A B\to\Omega^1_B$
deducido de $i$ (20.5.2) está, por tanto, en el sub-$B$-módulo de
$\Omega^1_B$ generado por los $d_B(i(j(b)))$ para $b\in B$; en virtud de
(21.3.1.2), esta imagen es nula. En otras palabras, se tiene una sucesión de
homomorfismos
\begin{equation}
\tag{21.3.3.4}\label{0.21.3.3.4-es}
0\to\Xi_{B/A}\otimes_A B_{[i]}\to(\Omega^1_{B/A})^{(p)}
\to\Upsilon_{B/A}
\end{equation}
que no es necesariamente exacta, pero en la que la composición de dos
homomorfismos consecutivos es nula.
\end{env}

\begin{proposition}[21.3.4]
\label{0.21.3.4-es}
Si $B$ es un $A$-módulo plano (para $i$), la sucesión (21.3.3.4) es exacta.
\end{proposition}
\begin{proof}
Esto se deduce de la definición de la planitud. La proposición se aplica en
particular cuando $B^p\subset A\subset B$, siendo $i$ y $j$ respectivamente
la inyección canónica y $F_B$, y $B$ admite una $p$-base sobre $A$ (de modo
que $B$ es entonces un $A$-módulo libre). Pero incluso en este caso el núcleo
$\Xi_{B/A}$ no es necesariamente nulo. Sin embargo:
\end{proof}

\begin{proposition}[21.3.5]
\label{0.21.3.5-es}
Sean $B$ un anillo reducido, $A$ un subanillo de $B$ tal que
$B^p\subset A\subset B$. Supongamos que existe una $p$-base
$(x_\lambda)_{\lambda\in L}$ de $B$ sobre $A$ y una $p$-base
$(y_\mu)_{\mu\in M}$ de $A$ sobre $B^p$; entonces el homomorfismo canónico
(21.3.3.4)
\begin{equation}
\tag{21.3.5.1}\label{0.21.3.5.1-es}
(\Omega^1_{B/A})^{(p)}\longrightarrow\Upsilon_{B/A}
\end{equation}
es biyectivo, y los elementos $d_A(x_\lambda^p)\otimes1$ forman una base del
$B$-módulo $\Upsilon_{B/A}$.
\end{proposition}

\begin{proof}
\oldpage[0\textsubscript{IV-1}]{162}
Como $B$ es reducido, $x\mapsto x^p$ es un isomorfismo estructural de $B$
sobre $B^p$; transportando la estructura mediante este isomorfismo, se ve
que $(x_\lambda^p)_{\lambda\in L}$ y los $y_\mu$ forman una $p$-base de
$B^p$ sobre $A^p$ (21.1.10); se deduce entonces de (21.1.5), (ii), y (21.2.5)
que los $d_A(x_\lambda^p)$ y los $d_A(y_\mu)$ forman una base del $A$-módulo
$\Omega^1_A$. Ahora bien, la imagen de $d_A(x_\lambda^p)\otimes1$ por
$i_{B/A}$ es $d_B(x_\lambda^p)=0$; por otra parte, la imagen de $d_A(y_\mu)$
por $i_{B/A}$ es $d_B(i(y_\mu))$, y como los $y_\mu$ y los $x_\lambda$
forman una $p$-base de $B$ sobre $B^p$ (21.1.10), los $d_B(y_\mu)$ son una
parte de una base de $\Omega^1_B$ (21.2.5), por lo que son linealmente
independientes sobre $B$. Se deduce que el núcleo de $i_{B/A}$ tiene por base
los $d_A(x_\lambda^p)\otimes1$; como estos son las imágenes por (21.3.5.1) de
los elementos correspondientes de $(\Omega^1_{B/A})^{(p)}$, y estos últimos
forman una base del $B$-módulo $(\Omega^1_{B/A})^{(p)}$ (21.2.5), el
homomorfismo (21.3.5.1) de $B$-módulos es biyectivo.
\end{proof}

\begin{remarks}[21.3.6]
\label{0.21.3.6-es}
\begin{enumerate}
  \item[(i)] Sean $A'$, $B'$ dos anillos de característica $p$, y supongamos
  dados dos homomorfismos $i':A'\to B'$, $j':B'\to A'$ que satisfacen
  (21.3.1.1) y (21.3.1.2), así como homomorfismos de anillos
  $f:A\to A'$, $g:B\to B'$ que hacen conmutativo el diagrama
  \[
  \xymatrix{
  A'\ar[r]^{i'}&B'\ar[r]^{j'}&A'\\
  A\ar[u]^f\ar[r]^i&B\ar[u]^g\ar[r]^j&A\ar[u]^f
  }
  \]
  Entonces el homomorfismo canónico $\Omega^1_{B/A}\to\Omega^1_{B'/A'}$
  (20.5.4.3) da un dihomomorfismo canónico
  $\Theta_{B/A}\to\Theta_{B'/A'}$ que hace conmutativo el diagrama
  \[
  \xymatrix{
  0\ar[r]&\Xi_{B/A}\ar[r]\ar[d]&\Theta_{B/A}\ar[r]^{\pi_{B/A}}\ar[d]&\Omega^1_A\ar[d]\\
  0\ar[r]&\Xi_{B'/A'}\ar[r]&\Theta_{B'/A'}\ar[r]^{\pi_{B'/A'}}&\Omega^1_{A'}
  }
  \]
  \item[(ii)] Sea $A'$ una $A$-álgebra cualquiera y pongamos
  $B'=B\otimes_A A'$; entonces $i'=i\otimes1:A'\to B'$ y
  $j'=j\otimes1:B'\to A'$ satisfacen (21.3.1.1) y (21.3.1.2), y se tiene
  $\Omega^1_{B'/A'}=\Omega^1_{B/A}\otimes_A A'=\Omega^1_{B/A}\otimes_B B'$
  (20.5.5); se deduce un $A'$-isomorfismo canónico
  \begin{equation}
  \tag{21.3.6.1}\label{0.21.3.6.1-es}
  \Theta_{B/A}\otimes_A A'\xrightarrow{\sim}\Theta_{B'/A'}
  \end{equation}
  y también un $B'$-homomorfismo canónico
  \begin{equation}
  \tag{21.3.6.2}\label{0.21.3.6.2-es}
  (\Omega^1_{B/A})^{(p)}\otimes_B B'\xrightarrow{\sim}
  (\Omega^1_{B'/A'})^{(p)}.
  \end{equation}
\end{enumerate}
\end{remarks}

\subsection{Caso de las extensiones de cuerpos}
\label{0.21.4-es}

\begin{env}[21.4.0]
\label{0.21.4.0-es}
Sean $K$ un cuerpo de característica $p>0$, $k$ un subcuerpo de $K$; entonces
el anillo $k[K^p]$ es igual al cuerpo $k(K^p)$, puesto que $k$ es algebraico
sobre $K^p$. Por tanto se pueden aplicar los resultados de los números
anteriores reemplazando en todas partes $A$, $B$ y $A[B^p]$ por $k$, $K$ y
$k(K^p)$.
\end{env}

\begin{lemma}[21.4.1]
\label{0.21.4.1-es}
Sean $k$ un cuerpo de característica $p>0$, $K$ una extensión de $k$. Para
que un elemento $x\in K$ sea $p$-libre sobre $k$, es necesario y suficiente
que $x\notin k(K^p)$.
\end{lemma}

\begin{proof}
En efecto, $x$ es raíz del polinomio $X^p-x^p$ de $k(K^p)[X]$, y se sabe
(Bourbaki, \emph{Alg.}, cap.~V, §~8, n$^\circ$~1, prop.~1) que, si
$x\notin k(K^p)$, este polinomio es irreducible, de modo que los elementos
$1,x,\ldots,x^{p-1}$ forman una base del $k(K^p)$-módulo $k(K^p)(x)$.
\end{proof}

\begin{theorem}[21.4.2]
\label{0.21.4.2-es}
Sean $k$ un cuerpo de característica $p>0$, $K$ una extensión de $k$,
$S$ un sistema de $p$-generadores de $K$ sobre $k$, $L\subset S$ una parte
$p$-libre sobre $k$. Existe entonces una $p$-base $B$ de $K$ sobre $k$ tal que
$L\subset B\subset S$. En particular, toda extensión de $k$ admite una
$p$-base sobre $k$.
\end{theorem}

\begin{proof}
Podemos limitarnos al caso en que $K^p\subset k$. En virtud del teorema de
Zorn, existe en $K$ una parte $B$ tal que $L\subset B\subset S$, $p$-libre
sobre $k$ y maximal entre las partes de $S$ con estas propiedades. Basta ver
que el subcuerpo $K'$ de $K$ generado por $k$ y $B$ es igual a $K$.
En caso contrario existiría $x\in S$ que no pertenece a $K'$; como
$K^p\subset k$, se tiene $K'=K'(K^p)$; por tanto $x$ sería $p$-libre sobre
$K'$ (21.4.1), y, en consecuencia, $B\cup\{x\}$ sería $p$-libre sobre $k$
(21.1.10), en contradicción con la definición de $B$. La última afirmación de
(21.4.2) se obtiene tomando $L=\varnothing$, $S=K$.
\end{proof}

\begin{corollary}[21.4.3]
\label{0.21.4.3-es}
Sean $k$ un cuerpo de característica $p>0$, $K$ una extensión de $k$.
Para que una familia $(x_\alpha)$ de elementos de $K$ sea $p$-independiente sobre $k$,
es necesario y suficiente que, para todo $\alpha$, $x_\alpha$ no pertenezca al
cuerpo $K_\alpha$ generado por $k(K^p)$ y por los $x_\beta$ de índice
$\beta\neq\alpha$.
\end{corollary}

\begin{proof}
La condición es necesaria en virtud de (21.1.10). Recíprocamente, supongamos que se
cumple; podemos limitarnos al caso en que $(x_\alpha)$ sea un sistema de
$p$-generadores de $K$. Existe entonces una subfamilia de $(x_\alpha)$ que es una
$p$-base de $K$ (21.4.2), pero esta familia no puede ser distinta de $(x_\alpha)$,
pues, en virtud de la hipótesis, no sería un sistema de $p$-generadores de $K$.
\end{proof}

\oldpage[0\textsubscript{IV-1}]{163}
\begin{corollary}[21.4.4]
\label{0.21.4.4-es}
Sean $k$ un cuerpo de característica $p>0$, $K$ una extensión de $k$.
Para que $\Omega^1_{K/k}=0$, es necesario y suficiente que $K=k(K^p)$.
En particular, para que $\Omega^1_K=0$, es necesario y suficiente que $K$ sea
un cuerpo perfecto.
\end{corollary}

\begin{proof}
En efecto, si $(x_\alpha)$ es una $p$-base de $K$ sobre $k$, los
$d_{K/k}(x_\alpha)$ forman una base del espacio vectorial $\Omega^1_{K/k}$
(21.2.5).
\end{proof}

\begin{theorem}[21.4.5]
\label{0.21.4.5-es}
Sean $k$ un cuerpo de característica $p>0$, $K$ una extensión de $k$,
$(x_\alpha)$ una familia de elementos de $K$. Para que $(x_\alpha)$ sea
$p$-independiente sobre $k$ (resp. una familia de $p$-generadores de $K$ sobre $k$,
resp. una $p$-base de $K$ sobre $k$), es necesario y suficiente que la familia
$(d_{K/k}(x_\alpha))$ sea una familia libre (resp. un sistema de
generadores, resp. una base) en el espacio vectorial $\Omega^1_{K/k}$.
\end{theorem}

\begin{proof}
\oldpage[0\textsubscript{IV-1}]{164}
Sea $K'$ el subcuerpo (igual al subanillo) de $K$ generado por
$k(K^p)$ y los $x_\alpha$; teniendo en cuenta (21.4.3), basta ver que,
para cada $\alpha$, $x_\alpha$ no pertenece a $K_\alpha$. Ahora bien, en la
sucesión exacta
\[
\Omega^1_{K_\alpha/k(K^p)}\otimes_{K_\alpha}K\longrightarrow
\Omega^1_{K/k(K^p)}\longrightarrow\Omega^1_{K/K_\alpha}\longrightarrow0,
\]
las imágenes por la flecha de la izquierda de los $d_{K_\alpha/k}(x_\beta)$ para
$\beta\neq\alpha$ son los $d_{K/k}(x_\beta)$, por lo que generan un subespacio
vectorial de $\Omega^1_{K/k(K^p)}$ que no contiene a $d_{K/k}(x_\alpha)$; como
este subespacio es el núcleo de la flecha de la derecha, se ve que
$d_{K/K_\alpha}(x_\alpha)\neq0$, y por tanto $x_\alpha\notin K_\alpha$.

Si ahora $(x_\alpha)$ es una familia $p$-independiente, es una subfamilia
de una $p$-base de $K$ sobre $k$ (21.4.2); los $d_{K/k}(x_\alpha)$ forman entonces
parte de una base del espacio vectorial $\Omega^1_{K/k}$ (21.2.5), y son
por consiguiente linealmente independientes sobre $K$. Recíprocamente, si los
$d_{K/k}(x_\alpha)$ son linealmente independientes sobre $K$, los dos
párrafos precedentes muestran que $(x_\alpha)$ es $p$-independiente sobre $k$.
\end{proof}

\begin{corollary}[21.4.6]
\label{0.21.4.6-es}
Sean $k$ un cuerpo de característica $p>0$, $K$ una extensión de $k$,
$x$ un elemento de $K$. Las tres condiciones siguientes son equivalentes:
\begin{enumerate}
  \item[a)] $x\notin k(K^p)$.
  \item[b)] $d_{K/k}(x)\neq0$.
  \item[c)] El elemento $x$ es $p$-independiente sobre $k$.
\end{enumerate}
\end{corollary}

\begin{proposition}[21.4.7]
\label{0.21.4.7-es}
Sean $L$ un cuerpo de característica $p>0$, $K$ un subcuerpo de $L$ tal
que $L^p\subset K\subset L$. Entonces la sucesión de espacios vectoriales sobre $L$
\begin{equation}
\tag{21.4.7.1}\label{0.21.4.7.1-es}
0\to\Omega^1_{K/L^p}\otimes_KL\to\Omega^1_{L/L^p}
\to\Omega^1_{L/K}\to0
\end{equation}
es exacta; en otros términos, se tiene $\Upsilon_{L/K/L^p}=0$.
\end{proposition}
\begin{proof}
Es un caso particular de (21.2.4), teniendo en cuenta (21.4.2).
\end{proof}

\begin{proposition}[21.4.8]
\label{0.21.4.8-es}
Bajo las hipótesis de (21.4.7), el homomorfismo canónico
\begin{equation}
\tag{21.4.8.1}\label{0.21.4.8.1-es}
(\Omega^1_{L/K})^{(p)}\longrightarrow\Upsilon_{L/K}
\end{equation}
es biyectivo; si $(x_\lambda)$ es una $p$-base de $L$ sobre $K$, los elementos
$d_K(x_\lambda^p)\otimes1$ forman una base del espacio vectorial $\Upsilon_{L/K}$.
\end{proposition}
\begin{proof}
Es un caso particular de (21.3.5).
\end{proof}

\subsection{Aplicación: criterios de separabilidad}
\label{0.21.5-es}

En este número y en los dos siguientes, ya no suponemos que los anillos
considerados sean de característica $p>0$.

\begin{env}[21.5.1]
\label{0.21.5.1-es}
Notemos primero que el criterio (21.2.7) permite demostrar una parte del
teorema de Cohen sobre las extensiones separables (19.6.1), a saber, que si $k$
es de característica $p>0$ y si $K$ es una extensión separable de $k$,
entonces $K$ es una $k$-álgebra formalmente lisa. En efecto, $K$ admite una
$p$-base sobre $k$ (21.4.2), y por otra parte se deduce del criterio de MacLane
(Bourbaki, \emph{Alg.}, cap.~V, §~8, n$^\circ$~2, prop.~3) que, en una
clausura algebraica de $K$, $k^{p^{-1}}$ y $K$ son linealmente disjuntos sobre
$k$, y por consiguiente que el homomorfismo canónico
$k^{p^{-1}}\otimes_k K\to k^{p^{-1}}(K)$ es biyectivo, lo que es precisamente
la condición (i) de (21.2.7), después de transportar la estructura mediante el
isomorfismo $x\mapsto x^p$.
\end{env}
\begin{proposition}[21.5.2]
\label{0.21.5.2-es}
\textnormal{(MacLane).}
Sean $A$, $B$ dos anillos de valoración discreta completos, $\mathfrak m$,
$\mathfrak n$ sus ideales máximos respectivos, $K=A/\mathfrak m$ y
$L=B/\mathfrak n$ sus cuerpos residuales, $u:A\to B$ un homomorfismo tal que
$Bu(\mathfrak m)=\mathfrak n$, $u_0=u\otimes 1:K\to L$ el homomorfismo
correspondiente para los cuerpos residuales. Consideremos las condiciones siguientes:
\begin{enumerate}
  \item[\rm a)] $L$ es una extensión separable de $K$ (para $u_0$).
  \item[\rm b)] Para todo anillo de valoración discreta completo $B'$ de ideal
  máximo $\mathfrak n'$ y de cuerpo residual $L'$,
\oldpage[0\textsubscript{IV-1}]{165}
  todo homomorfismo
  $u':A\to B'$ tal que $B'u'(\mathfrak m)=\mathfrak n'$ y todo
  $K$-isomorfismo $\sigma:L\to L'$ (relativo a $u_0$ y
  $u'_0=u'\otimes1:K\to L'$), existe un isomorfismo $w:B\to B'$ tal que
  $u'=w\circ u$ y $w_0=w\otimes1:L\to L'$ sea igual a $\sigma$.
  \item[\rm b$'$)] Para todo homomorfismo $u':A\to B$ tal que
  $Bu'(\mathfrak m)=\mathfrak n$ y tal que $u'_0=u'\otimes1:K\to L$ sea igual
  a $u_0$, existe un automorfismo $w$ de $B$ tal que $u'=w\circ u$ y
  $w_0=w\otimes1:L\to L$ sea la identidad.
  \item[\rm c)] Designando por $u_1:A/\mathfrak m^2\to B/\mathfrak n^2$
  el homomorfismo deducido de $u$ por paso a los cocientes, entonces, para todo
  homomorfismo local $u'_1:A/\mathfrak m^2\to B/\mathfrak n^2$ tal que
  $\operatorname{gr}_0(u'_1)=\operatorname{gr}_0(u_1)(=u_0)$, existe un
  automorfismo $w_1$ de $B/\mathfrak n^2$ tal que
  $\operatorname{gr}_0(w_1)$ sea la identidad y $u'_1=w_1\circ u_1$.
  \item[\rm c$'$)] Para todo homomorfismo local $u'_1:A/\mathfrak m^2\to
  B/\mathfrak n^2$ tal que $\operatorname{gr}_0(u'_1)=\operatorname{gr}_0(u_1)$
  y $\operatorname{gr}_1(u'_1)=\operatorname{gr}_1(u_1)$
  (homomorfismo de $\mathfrak m/\mathfrak m^2$ en
  $\mathfrak n/\mathfrak n^2$), existe un automorfismo $w_1$ de $B/\mathfrak n^2$
  tal que $\operatorname{gr}_0(w_1)$ sea la identidad y $u'_1=w_1\circ u_1$.
\end{enumerate}
Entonces se tienen las implicaciones $c)\Leftrightarrow c')\Leftrightarrow a)\Rightarrow b)
 \Rightarrow b')$. Si además $A$ es un anillo de Cohen, las cinco condiciones
anteriores son equivalentes.

\end{proposition}
\begin{proof}
Las implicaciones $b)\Rightarrow b')$ y $c)\Rightarrow c')$ son triviales.
Demostremos que $a)$ implica $b)$.
El homomorfismo $u$ hace de $B$ un $A$-módulo sin torsión, pues transforma por
hipótesis un uniformizante de $A$ en un uniformizante de $B$; se deduce que
$B$ es un $A$-módulo plano (0$_{\rm I}$, 6.3.4), luego un $A$-álgebra de
Cohen (19.8.1); el hecho de que $a)$ implique $b)$ es entonces consecuencia de
(19.8.2, (i)) aplicada a $G=B'$, $\mathfrak J=\mathfrak n'$,
considerando $B'$ como $A$-álgebra mediante $u'$ y siendo el homomorfismo
$B\to B'/\mathfrak n'$ el compuesto
$B\to B/\mathfrak n\to B'/\mathfrak n'$, que es un $A$-homomorfismo en
virtud de la hipótesis $u'_0=\sigma\circ u_0$. El mismo razonamiento prueba que
$a)$ implica $c)$, tomando esta vez $G=B/\mathfrak n^2$,
$\mathfrak J=\mathfrak n/\mathfrak n^2$, y considerando $G$ como
$A$-álgebra mediante $u'_1$.
Probemos en segundo lugar que $c')$ implica $a)$. Los dos homomorfismos
$u_1$ y $u'_1$ son tales que $u'_1=u_1+D$, donde $D$ es una derivación de
$A/\mathfrak m^2$ en $\mathfrak n/\mathfrak n^2$ (20.1.1); además,
la hipótesis $\operatorname{gr}_1(u'_1)=\operatorname{gr}_1(u_1)$ significa
que $D$ es nula en $\mathfrak m/\mathfrak m^2$, y por tanto puede
considerarse como una derivación de $K=A/\mathfrak m$ en
$\mathfrak n/\mathfrak n^2$, y $\mathfrak n/\mathfrak n^2$ se identifica
(eligiendo un uniformizante de $B$) con el $K$-módulo $L$ (para $u_0$).
Por otra parte, las condiciones impuestas a $w_1$ implican que
$\operatorname{gr}_1(w_1)$ es también la identidad (pues
$u_1(\mathfrak m/\mathfrak m^2)$ genera
$\mathfrak n/\mathfrak n^2$); por tanto $w_1$ es de la forma
$x\mapsto x+D'(x)$, donde $D'$ es ahora una derivación de
$B/\mathfrak n^2$ en el $(B/\mathfrak n^2)$-módulo
$\mathfrak n/\mathfrak n^2$; como $w_1$ es la identidad en
$\mathfrak n/\mathfrak n^2$, $D'$ puede también (mediante la identificación
anterior de $\mathfrak n/\mathfrak n^2$ con $L$) considerarse como una
derivación de $L$ en $L$; finalmente, la relación $u'_1=w_1\circ u_1$
significa que $D'$ prolonga $D$. Por otra parte, toda derivación $D$
de $K$ en $L$ corresponde a un homomorfismo $u'_1$ que verifica las condiciones
de $c')$ (20.1.1); por tanto, la condición $c')$ significa que toda derivación de
$K$ en $L$ se prolonga a una derivación de $L$ en sí mismo, es decir,
(20.6.5), que $L$ es separable sobre $K$.
Probemos finalmente que, cuando $A$ es un anillo de Cohen, $b')$ implica $c')$.
Probemos primero que bajo las hipótesis de $c')$, existe un homomorfismo
$u':A\to B$ que verifica las hipótesis de $b')$ y que, por paso a los cocientes,
da $u'_1:A/\mathfrak m^2\to B/\mathfrak n^2$. En efecto, esto se deduce de
(19.8.6, (i)) aplicado al homomorfismo compuesto
$v':A\to A/\mathfrak m^2\xrightarrow{u'_1}B/\mathfrak n^2$;
\oldpage[0\textsubscript{IV-1}]{166}
este último se factoriza entonces como $A\xrightarrow{u'}B\to B/\mathfrak n^2$
y las hipótesis sobre $\operatorname{gr}_0(u'_1)$ y
$\operatorname{gr}_1(u'_1)$ implican
$\operatorname{gr}_0(u')=\operatorname{gr}_0(u)=u_0$ y
$\operatorname{gr}_1(u')=\operatorname{gr}_1(u)$, de modo que la imagen por
$u'$ de un uniformizante de $A$ es un uniformizante de $B$, y por tanto
$Bu'(\mathfrak m)=\mathfrak n$. Basta entonces, para obtener un automorfismo
$w_1$ que responda a la cuestión, tomar el automorfismo inducido por paso a los
cocientes por el automorfismo $w$ de $B$ proporcionado por la aplicación de
$b')$ al homomorfismo $u'$.
\end{proof}
\begin{remarks}[21.5.3]
\label{0.21.5.3-es}
\begin{enumerate}
  \item[(i)] Las propiedades diferenciales de los cuerpos permiten resolver la
  cuestión de la unicidad del cuerpo de representantes en un anillo local
  noetheriano completo $C$ (19.8.7, (ii)). Sean $\mathfrak J$ el ideal máximo
  de $C$, $K=C/\mathfrak J$ el cuerpo residual de $C$; podemos limitarnos al
  caso $\mathfrak J\ne0$. Supongamos que existe un homomorfismo $u:K\to C$
  que, compuesto con la augmentación $C\to K$, da la identidad; entonces, para
  que este homomorfismo sea único, es necesario y suficiente que
  $\Omega^1_K=0$. En efecto, la condición $\Omega^1_K=0$ implica que $K$ es
  formalmente no ramificado sobre su cuerpo primo (20.7.4); si existiera un
  segundo homomorfismo $v\ne u$ que respondiera a la cuestión, habría un
  entero $n$ máximo tal que $\mathfrak J^n$ contuviese el conjunto de
  $u(x)-v(x)$ para $x\in K$; por paso al cociente, $u$ y $v$ darían dos
  homomorfismos distintos $u_1,v_1$ de $K$ en
  $C/\mathfrak J^{n+1}$, cuyos compuestos con
  $C/\mathfrak J^{n+1}\to C/\mathfrak J^n$ serían iguales, lo que contradice
  la definición (19.10.2). Recíprocamente, supongamos que $\Omega^1_K\ne0$;
  existe entonces una derivación $D\ne0$ de $K$ en
  $\mathfrak J/\mathfrak J^2$ (20.4.8), y por tanto un homomorfismo
  $v_1:K\to C/\mathfrak J^2$ tal que, si $u_1:K\to C/\mathfrak J^2$ se
  obtiene de $u$ por paso al cociente, se tiene $v_1=u_1+D$
  (20.1.1). Si $k$ es el cuerpo primo de $K$, $C$ es una $k$-álgebra y
  $K$ una $k$-álgebra formalmente lisa, y las restricciones de
  $u_1$ y $v_1$ a $k$ coinciden, por lo que $v_1$ se factoriza como
  $K\xrightarrow{v}C\to C/\mathfrak J^2$ y se tiene $v\ne u$.
  Recordemos ((20.6.20) y (21.4.4)) que la condición $\Omega^1_K=0$ significa
  que $K$ es perfecto si es de característica $\ne0$, y una extensión
  algebraica de $\mathbf Q$ si es de característica $0$.
  \item[(ii)] De la misma manera, sean $W$ un anillo de Cohen, $C$ un anillo
  local noetheriano completo, y sea $\mathfrak J\ne0$ un ideal de $C$ contenido en
  el ideal máximo; para que la factorización $W\to C\to C/\mathfrak J$
  en (19.8.6, (i)) sea única, es necesario y suficiente que el cuerpo residual
  $K$ del anillo de Cohen $W$ satisfaga $\Omega^1_K=0$. En efecto, si
  $\Omega^1_K\ne0$ y si $\mathfrak m$ designa el ideal máximo de $W$, basta
  componer una derivación $D\ne0$ de $K$ en
  $\mathfrak J/\mathfrak m\mathfrak J$ con la augmentación $W\to K$ para
  obtener una derivación $D'\ne0$ de $W$ en
  $\mathfrak J/\mathfrak m\mathfrak J$ y se termina el razonamiento como
  en (i), formando mediante $D'$ un homomorfismo
  $v_1:W\to C/\mathfrak m\mathfrak J$ distinto del homomorfismo
  $u_1:W\to C/\mathfrak m\mathfrak J$ obtenido de $u$ por paso al cociente;
  se concluye invocando esta vez (19.8.6, (i)). Si, por el contrario,
  $\Omega^1_K=0$, la unicidad de $v$ ya resulta de (i) cuando $W=K$ es un
  cuerpo de característica $0$. En el caso contrario, se tiene
  $\widehat{\Omega}^1_W=0$; en efecto, entonces $\mathfrak m=pW$, siendo $p$
  la característica de $K$ (19.8.5), y el homomorfismo canónico
  $\mathfrak m/\mathfrak m^2\to\Omega^1_W/\mathfrak m\Omega^1_W$ (20.5.11.2)
  es por tanto nulo. La sucesión exacta (20.5.12.1) aplicada a $W$ y a
  $K=W/\mathfrak m$ implica entonces que $\Omega^1_W=\mathfrak m\Omega^1_W$,
  de donde se sigue nuestra afirmación. Pero entonces (20.7.4) $W$ es
  formalmente no ramificado (para su topología adítica) sobre $\mathbf Z$, y la
  unicidad de $v$ se demuestra como en (i).
\end{enumerate}
\end{remarks}

\oldpage[0\textsubscript{IV-1}]{167}
\subsection{Cuerpos admisibles para una extensión}
\label{section:0.21.6-es}

\begin{env}[21.6.1]
\label{0.21.6.1-es}
Dados cuatro cuerpos $k_0\subset k\subset K\subset L$, se deduce de
(20.6.16) y (20.6.17) que se tiene una sucesión exacta
\begin{equation}
\tag{21.6.1.1}
\label{0.21.6.1.1-es}
0\to\Upsilon_{K/k/k_0}\otimes_KL\xrightarrow{v}\Upsilon_{L/k/k_0}
\xrightarrow{u}\Upsilon_{L/K/k_0}\xrightarrow{s}\Upsilon_{L/K/k}\to0.
\end{equation}

Cuando fijamos $k_0$, $K$ y $L$ y hacemos «variar» el cuerpo intermedio
$k$ entre $k_0$ y $K$, evidentemente se tiene
$\Upsilon_{L/K/k_0}=\Upsilon_{L/K/k}$ cuando $k=k_0$. Cuando el homomorfismo
canónico $s$ de (21.6.1.1) es además biyectivo, decimos que $k$ es un cuerpo
$k_0$-admisible para la extensión $L$ de $K$. El interés de la existencia, bajo
ciertas condiciones, de tales cuerpos $k$, que sean sin embargo
«suficientemente cercanos» a $K$ (por ejemplo, tales que $[K:k]$ sea finito), es
que permiten sustituir en ciertas cuestiones los módulos de diferenciales
$\Omega^1_{K/k_0}$ y $\Omega^1_{L/k_0}$ (que pueden ser «demasiado grandes»,
por ejemplo cuando $k_0$ es el cuerpo primo) por
$\Omega^1_{K/k}$ y $\Omega^1_{L/k}$, más fáciles de manejar.

Cuando $k_0$ es el \emph{cuerpo primo}, diremos «cuerpo admisible» en lugar
de «cuerpo $k_0$-admisible».
 
Se pondrá
\end{env}

\begin{env}[21.6.1.2]
\label{0.21.6.1.2-es}
\[
\Delta(L/K,k/k_0)=\operatorname{Coker}
(\Upsilon_{K/k/k_0}\otimes_KL\to\Upsilon_{L/k/k_0})
\simeq\operatorname{Ker}(\Upsilon_{L/K/k_0}\to\Upsilon_{L/K/k})
\]
(espacio vectorial sobre $L$); su rango se denotará $d(L/K,k/k_0)$ y se llamará
\emph{defecto de $k_0$-admisibilidad de $k$ para la extensión $L$ de $K$}
(es evidentemente nulo si y solo si $k$ es $k_0$-admisible para esta
extensión). Cuando $k_0$ es el cuerpo primo, escribiremos
$\Delta(L/K,k)$ y $d(L/K,k)$ en lugar de $\Delta(L/K,k/k_0)$ y
$d(L/K,k/k_0)$.
\end{env}

\begin{proposition}[21.6.2]
\label{0.21.6.2-es}
Sean $k_0\subset k\subset K\subset L$ cuatro cuerpos.
\begin{enumerate}
  \item[\rm (i)] Las condiciones siguientes son equivalentes:
  \begin{enumerate}
    \item[\rm a)] El cuerpo $k$ es $k_0$-admisible para la extensión $L$ de $K$
    (en otras palabras, el homomorfismo $\Upsilon_{L/K/k_0}\to
    \Upsilon_{L/K/k}$ es inyectivo y, por tanto, biyectivo).
    \item[\rm b)] El homomorfismo canónico
    $u:\Upsilon_{K/k/k_0}\to\Upsilon_{L/k/k_0}$ es nulo.
    \item[\rm c)] El homomorfismo canónico
    $v:\Upsilon_{K/k/k_0}\otimes_KL\to\Upsilon_{L/k/k_0}$ es sobreyectivo
    (y, por tanto, biyectivo).
    \item[\rm d)] Se tiene $d(L/K,k/k_0)=0$ (o
    $\Delta(L/K,k/k_0)=0$).
  \end{enumerate}
  \item[\rm (ii)] Las condiciones equivalentes de (i) se verifican cuando
  se está en uno de los casos siguientes: $\alpha)$ $L$ es separable sobre $k$;
  $\beta)$ $L$ es separable sobre $K$; $\gamma)$ se tiene
  $k\subset k_0(K^p)$, designando por $p$ el exponente de característica de $k_0$.
\end{enumerate}
\end{proposition}

\begin{proof}
\begin{enumerate}
  \item[(i)] Las afirmaciones se deducen trivialmente de la exactitud de la sucesión
  (21.6.1.1).
  \item[(ii)] Si $L$ es separable sobre $k$, se tiene
  $\Upsilon_{L/k/k_0}=0$ (20.6.19), por lo que se cumple la condición b) de (i);
  si $L$ es separable sobre $K$, se tiene
  $\Upsilon_{L/K/k_0}=0$ (20.6.19), por lo que se cumple la condición a) de (i);
  finalmente, si se tiene $k\subset k_0(K^p)$, se deduce que
  $\Upsilon_{L/K/k_0}=\Upsilon_{L/K/k}$ en virtud de (21.1.5.1), y se cumple
  la condición a) de (i).
\end{enumerate}
\end{proof}
\oldpage[0\textsubscript{IV-1}]{168}
\begin{env}[21.6.3]
\label{0.21.6.3-es}
Supongamos que tenemos un diagrama conmutativo de monomorfismos de cuerpos
\[
\xymatrix{
 k'_0\ar[r]&k'\ar[r]&K'\ar[r]&L'\\
 k_0\ar[u]\ar[r]&k\ar[u]\ar[r]&K\ar[u]\ar[r]&L\ar[u]
}
\]
Se deduce entonces de (20.6.17, (ii)) que se tiene un homomorfismo canónico
\begin{equation}
\tag{21.6.3.1}
\label{0.21.6.3.1-es}
\Delta(L/K,k/k_0)\longrightarrow\Delta(L'/K',k'/k'_0)
\end{equation}
con una propiedad de transitividad evidente, de modo que podemos decir que
$\Delta(L/K,k/k_0)$ es un \emph{functor} en el cuádruple
$(k_0,k,K,L)$.
\end{env}

\begin{proposition}[21.6.4]
\label{0.21.6.4-es}
\begin{enumerate}
  \item[\rm (i)] Sean $k\subset k'\subset k''\subset K\subset L$ cinco cuerpos.
  Se tiene una sucesión exacta de homomorfismos canónicos
\begin{equation}
\tag{21.6.4.1}
\label{0.21.6.4.1-es}
0\to\Delta(L/K,k'/k)\to\Delta(L/K,k''/k)\to\Delta(L/K,k''/k')\to0
\end{equation}
y, por consiguiente, la igualdad
\begin{equation}
\tag{21.6.4.2}
\label{0.21.6.4.2-es}
d(L/K,k''/k)=d(L/K,k''/k')+d(L/K,k'/k).
\end{equation}
  \item[\rm (ii)] Sean $k_0\subset k\subset K\subset L\subset M$ cinco cuerpos.
  Se tiene una sucesión exacta de homomorfismos canónicos
\begin{equation}
\tag{21.6.4.3}
\label{0.21.6.4.3-es}
0\to\Delta(L/K,k/k_0)\otimes_LM\to\Delta(M/K,k/k_0)
\to\Delta(M/L,k/k_0)\to0
\end{equation}
y, por consiguiente, la igualdad
\begin{equation}
\tag{21.6.4.4}
\label{0.21.6.4.4-es}
d(M/K,k/k_0)=d(M/L,k/k_0)+d(L/K,k/k_0).
\end{equation}
\end{enumerate}
\end{proposition}

\begin{proof}
\begin{enumerate}
  \item[(i)] Consideremos el diagrama conmutativo
\[
\xymatrix{
0\ar[r]&\Upsilon_{K/k'/k}\otimes_KL\ar[r]\ar[d]&
\Upsilon_{K/k''/k}\otimes_KL\ar[r]\ar[d]&
\Upsilon_{K/k''/k'}\otimes_KL\ar[r]\ar[d]&0\\
0\ar[r]&\Upsilon_{L/k'/k}\ar[r]\ar[d]&
\Upsilon_{L/k''/k}\ar[r]\ar[d]&
\Upsilon_{L/k''/k'}\ar[r]\ar[d]&0\\
0\ar[r]&\Delta(L/K,k'/k)\ar[r]&\Delta(L/K,k''/k)\ar[r]&
\Delta(L/K,k''/k')\ar[r]&0
}
\]
donde podemos considerar las tres filas como \emph{complejos}
$T_1,T_2,T_3$, respectivamente; la sucesión exacta (21.6.1.1) y la definición
de $\Delta$ (21.6.1.2) muestran que se tiene una sucesión exacta de complejos
$0\to T_1\to T_2\to T_3\to0$; apliquémosle la sucesión exacta de
\oldpage[0\textsubscript{IV-1}]{169}
cohomología, y observemos que, en virtud de la exactitud de (21.6.1.1), la
cohomología de $T_1$ y la de $T_2$ son nulas salvo en un mismo grado, para el
cual los módulos de cohomología son ambos iguales a
$\Upsilon_{k''/k'/k}\otimes_{k''}L$; como $T_1$ y $T_2$ tienen así la misma
cohomología, la de $T_3$ es necesariamente nula, lo que prueba (i).
  \item[(ii)] Consideremos del mismo modo el diagrama conmutativo
\[
\xymatrix{
0\ar[r]&\Delta(L/K,k/k_0)\otimes_LM\ar[r]\ar[d]&
\Delta(M/K,k/k_0)\ar[r]\ar[d]&\Delta(M/L,k/k_0)\ar[r]&0\\
0\ar[r]&\Upsilon_{L/K/k_0}\otimes_LM\ar[r]\ar[d]&
\Upsilon_{M/K/k_0}\ar[r]\ar[d]&\Upsilon_{M/L/k_0}\ar[r]&0\\
0\ar[r]&\Upsilon_{L/K/k}\otimes_LM\ar[r]&
\Upsilon_{M/K/k}\ar[r]&\Upsilon_{M/L/k}\ar[r]&0
}
\]
donde de nuevo consideramos las tres filas como \emph{complejos}
$T'_1,T'_2,T'_3$; la sucesión exacta (21.6.1.1) y la definición de $\Delta$
(21.6.1.2) dan aquí una sucesión exacta de complejos
$0\to T'_1\to T'_2\to T'_3\to0$, a la que aplicamos nuevamente la sucesión
exacta de cohomología; esta vez, en virtud de la exactitud de (21.6.1.1), la
cohomología de $T'_2$ y la de $T'_3$ son nulas salvo en un mismo grado, para el
cual los módulos de cohomología son ambos iguales a
$\Upsilon_{M/L/k_0}$; concluimos por tanto aquí que la cohomología de $T'_1$ es
nula, lo que establece (ii).
\end{enumerate}
\end{proof}

\begin{corollary}[21.6.5]
\label{0.21.6.5-es}
\begin{enumerate}
  \item[(i)] Dados cinco cuerpos $k\subset k'\subset k''\subset K\subset L$,
  para que $k''$ sea $k$-admisible para la extensión $L$ de $K$, es necesario y
  suficiente que $k'$ sea $k$-admisible y que $k''$ sea $k'$-admisible
  para la extensión $L$ de $K$.
  \item[(ii)] Dados cinco cuerpos
  $k_0\subset k\subset K\subset L\subset M$, para que $k$ sea
  $k_0$-admisible para la extensión $M$ de $K$, es necesario y suficiente que lo
  sea para la extensión $L$ de $K$ y para la extensión $M$ de $L$.
\end{enumerate}

Esto se deduce inmediatamente de las relaciones (21.6.4.2) y (21.6.4.4), siendo
los valores de $d$ $\geqslant0$.
\end{corollary}

\begin{corollary}[21.6.6]
\label{0.21.6.6-es}
Sean $k_0\subset k\subset K\subset L$ cuatro cuerpos, y supongamos que $k$
sea $k_0$-admisible para la extensión $L$ de $K$. Entonces, si $k'_0$, $k'$,
$K'$, $L'$ son cuatro cuerpos tales que
$k_0\subset k'_0\subset k'\subset k\subset K\subset K'\subset L'\subset L$,
$k'$ es $k'_0$-admisible para la extensión $L'$ de $K'$.
\end{corollary}

\subsection{La igualdad de Cartier}
\label{subsection:0.21.7-es}

El resultado siguiente traduce en términos de diferenciales un teorema de
MacLane sobre las derivaciones:

\oldpage[0\textsubscript{IV-1}]{170}
\begin{theorem}[21.7.1]
\label{0.21.7.1-es}
\textnormal{(Cartier).}
Sean $K$ un cuerpo, $L$ una extensión de tipo finito de $K$. Entonces
$\Omega^1_{L/K}$ y $\Upsilon_{L/K}$ son espacios vectoriales de rango
finito sobre $L$, y se tiene
\begin{equation}
\tag{21.7.1.1}
\label{0.21.7.1.1-es}
\operatorname{rg}\Omega^1_{L/K}-\operatorname{rg}\Upsilon_{L/K}
=\deg.\operatorname{tr}_K L.
\end{equation}
\end{theorem}

\begin{proof}
Si $L'$ es un cuerpo tal que $K\subset L'\subset L$, se tiene la sucesión exacta
(20.6.15.1) (aplicada a $\Lambda=P$, cuerpo primo de $K$, $A=K$,
$B=L'$, $C=L$)
\[
0\to\Upsilon_{L'/K}\otimes_{L'}L\to\Upsilon_{L/K}
\to\Omega^1_{L'/K}\otimes_{L'}L\to\Omega^1_{L/K}
\to\Omega^1_{L/L'}\to0
\]
de donde (0$_{\rm III}$, 11.10.2)
\[
\operatorname{rg}_L\Omega^1_{L/K}-\operatorname{rg}_L\Upsilon_{L/K}
=(\operatorname{rg}_{L'}\Omega^1_{L/L'}-\operatorname{rg}_{L'}\Upsilon_{L/L'})
+(\operatorname{rg}_{L'}\Omega^1_{L'/K}-\operatorname{rg}_{L'}\Upsilon_{L'/K}).
\]
Como, por otra parte, $\deg.\operatorname{tr}_K L=\deg.\operatorname{tr}_K L'
+\deg.\operatorname{tr}_{L'}L$, se ve, por inducción sobre el número de
generadores de la extensión $L$, que basta demostrar (21.7.1.1) cuando
$L=K(x)$. Distingamos entonces tres casos:
\begin{enumerate}
  \item[\rm a)] $x$ es trascendente sobre $K$; como $L$ es separable sobre
  $K$, se tiene $\Upsilon_{L/K}=0$ (20.6.3); por otra parte, (20.5.9) y
  (20.4.13) muestran que $\Omega^1_{L/K}$ tiene rango $1$, de donde se
  obtiene (21.7.1.1) en este caso.
  \item[\rm b)] $L$ es una extensión algebraica separable de $K$, de modo que
  se tiene de nuevo $\Upsilon_{L/K}=0$ (20.6.3). Por otra parte,
  $\Omega^1_{L/K}=0$ por (20.6.20), y se obtiene de nuevo (21.7.1.1).
  \item[\rm c)] $L$ es una extensión algebraica inseparable de $K$; el
  razonamiento inicial muestra que podemos limitarnos al caso $L^p\subset K$;
  se deduce entonces de (21.4.8) que
  $\operatorname{rg}_L\Upsilon_{L/K}=\operatorname{rg}_L\Omega^1_{L/K}$,
  de donde se obtiene de nuevo (21.7.1.1).
\end{enumerate}
\end{proof}

\begin{corollary}[21.7.2]
\label{0.21.7.2-es}
Sean $K$ un cuerpo, $L$ una extensión de tipo finito de $K$, $k$ un subcuerpo
de $K$. Entonces $\Omega^1_{L/K}$ y $\Upsilon_{L/K/k}$ son espacios
vectoriales de rango finito sobre $L$, y se tiene
\begin{equation}
\tag{21.7.2.1}
\label{0.21.7.2.1-es}
\operatorname{rg}\Omega^1_{L/K}-\operatorname{rg}\Upsilon_{L/K/k}
=\deg.\operatorname{tr}_K L+d(L/K,k).
\end{equation}
Por consiguiente, se tiene la desigualdad
\begin{equation}
\tag{21.7.2.2}
\label{0.21.7.2.2-es}
\operatorname{rg}\Omega^1_{L/K}-\operatorname{rg}\Upsilon_{L/K/k}
\geqslant\deg.\operatorname{tr}_K L.
\end{equation}
Además, para que los dos miembros de (21.7.2.2) sean iguales, es necesario y
suficiente que $k$ sea un cuerpo admisible para la extensión $L$ de $K$.
\end{corollary}

\begin{proof}
En efecto, como el homomorfismo
$\Upsilon_{L/K}\to\Upsilon_{L/K/k}$ es sobreyectivo, se tiene por definición
(21.6.1.2) $\operatorname{rg}\Upsilon_{L/K}-\operatorname{rg}\Upsilon_{L/K/k}
=d(L/K,k)$, y el corolario se deduce inmediatamente de (21.7.1) y de
(21.6.2).
\end{proof}

\begin{corollary}[21.7.3]
\label{0.21.7.3-es}
Sean $k\subset K\subset L$ tres cuerpos tales que $L$ sea una extensión de
tipo finito de $k$. Entonces se tiene
\begin{equation}
\tag{21.7.3.1}
\label{0.21.7.3.1-es}
\operatorname{rg}_L\Omega^1_{L/k}-\operatorname{rg}_K\Omega^1_{K/k}
=\deg.\operatorname{tr}_K L+d(L/K,k).
\end{equation}
\oldpage[0\textsubscript{IV-1}]{171}
Por consiguiente, se tiene la desigualdad
\begin{equation}
\tag{21.7.3.2}
\label{0.21.7.3.2-es}
\operatorname{rg}_L\Omega^1_{L/k}-\operatorname{rg}_K\Omega^1_{K/k}
\geqslant\deg.\operatorname{tr}_K L.
\end{equation}
La igualdad se alcanza si y solo si $k$ es un cuerpo admisible para
la extensión $L$ de $K$.
\end{corollary}

\begin{proof}
En efecto, se tiene la sucesión exacta (20.6.1.1)
\[
0\to\Upsilon_{L/K/k}\to\Omega^1_{K/k}\otimes_KL\to\Omega^1_{L/k}
\to\Omega^1_{L/K}\to0
\]
por lo que el primer miembro de (21.7.3.1) es igual a
\[
\operatorname{rg}_L\Omega^1_{L/k}-\operatorname{rg}_L\Upsilon_{L/K/k}
\]
y basta aplicar (21.7.2).
\end{proof}

El interés de este último corolario es que solo intervienen en él módulos de
diferenciales, con exclusión de los módulos de imperfección. Veremos además más
abajo (21.8.6) que para toda extensión de tipo finito $L$ de $K$ existe un
subcuerpo $k$ de $K$ tal que $[K:k]$ sea finito y que sea admisible para la
extensión $L$ de $K$, de modo que los dos miembros de (21.7.3.2) sean entonces
iguales.

\begin{corollary}[21.7.4]
\label{0.21.7.4-es}
Sea $K$ una extensión de tipo finito de un cuerpo $k$. Las condiciones siguientes
son equivalentes:
\begin{enumerate}
  \item[\rm a)] $K$ es una extensión finita separable de $k$.
  \item[\rm b)] $K$ es una $k$-álgebra formalmente no ramificada (19.10.2).
  \item[\rm c)] $K$ es una $k$-álgebra formalmente étale (19.10.2).
  \item[\rm d)] Se tiene $\Omega^1_{K/k}=0$.
\end{enumerate}
\end{corollary}

\begin{proof}
En efecto, se sabe (siendo discretas las topologías) que las condiciones b) y
d) son equivalentes (20.7.4) y que a) implica que $K$ es una $k$-álgebra
formalmente lisa (19.6.1), por lo que la conjunción de a) y b) equivale a c).
Todo se reduce a ver que d) implica a). Ahora bien, en virtud de (21.7.1.1),
la relación $\Omega^1_{K/k}=0$ implica que $K$ es algebraico (y por tanto finito)
sobre $k$ y que $\Upsilon_{K/k}=0$, es decir (20.6.3), que $K$ es separable
sobre $k$.
\end{proof}

\begin{remark}[21.7.5]
\label{0.21.7.5-es}
En virtud de (20.6.3.2) y de (0$_{\rm III}$, 11.10.2), el primer miembro de
(21.7.1.1) no es otro que la \emph{característica de Euler-Poincaré} del
complejo $K_\bullet(C/B/A)$ introducido en (20.6.3), para $C=L$, $B=K$ y
$A=P$ (cuerpo primo de $K$). En el capítulo dedicado a la teoría de
intersecciones y al teorema de Riemann--Roch, también desempeñarán un papel
importante las características de Euler-Poincaré generalizadas (con valores
en grupos de clases de $\mathcal O_X$-Módulos) de complejos que generalizan
los complejos $F_\bullet(C/A)$ considerados en (20.6.22).
\end{remark}

\subsection{Criterios de admisibilidad}
\label{subsection:0.21.8-es}

Volvemos a nuestras convenciones anteriores y suponemos, por tanto, que todos
los cuerpos considerados en este número son de característica $p>0$.

\begin{lemma}[21.8.1]
\label{0.21.8.1-es}
Sean $K$ un cuerpo, $k$ un subcuerpo de $K$, $(k_\alpha)_{\alpha\in I}$ una
familia filtrante decreciente de subcuerpos de $K$ tal que
$k=\bigcap_{\alpha\in I}k_\alpha$. Sean $V$ un espacio vectorial sobre $K$,
$(a_i)_{1\leqslant i\leqslant n}$ una familia finita de vectores de $V$; si la
familia $(a_i)$ es libre sobre $k$, existe un índice $\gamma$ tal que también
sea libre sobre $k_\gamma$.
\end{lemma}

\oldpage[0\textsubscript{IV-1}]{172}
\begin{proof}
Sea $r$ el rango de la familia $(a_i)_{1\leqslant i\leqslant n}$ sobre $K$, y
razonemos por inducción sobre $n-r$; la proposición es evidente para $n=r$,
pues entonces la familia $(a_i)_{1\leqslant i\leqslant r}$ es libre sobre $K$, y
por tanto sobre todo subcuerpo de $K$. Supongamos, por ejemplo, que
$(a_i)_{1\leqslant i\leqslant r}$ sea libre sobre $K$, y escribamos
$a_{r+1}=\sum_{i=1}^r\lambda_i a_i$ con $\lambda_i\in K$; como la familia
$(a_i)_{1\leqslant i\leqslant r+1}$ es libre sobre $k$, los $\lambda_i$ no
pueden pertenecer todos a $k$; supongamos, por ejemplo, que
$\lambda_{i_0}\notin k$. Entonces existe un índice $\beta$ tal que
$\lambda_{i_0}\notin k_\beta$; concluimos que la familia
$(a_i)_{1\leqslant i\leqslant r+1}$ es libre sobre $k_\beta$; en efecto, como la
familia $(a_i)_{1\leqslant i\leqslant r}$ es libre sobre todo subcuerpo de
$K$, si la familia $(a_i)_{1\leqslant i\leqslant r+1}$ no fuera libre sobre
$k_\beta$, $a_{r+1}$ sería igual a una combinación lineal de los $a_i$ con
coeficientes en $k_\beta$, y como estos coeficientes son necesariamente los
$\lambda_i$, se llega a una contradicción. Basta ahora aplicar la hipótesis de
inducción sustituyendo $K$ por $k_\beta$ y la familia $(k_\alpha)_{\alpha\in I}$
por la subfamilia de los $k_\alpha$ contenidos en $k_\beta$.
\end{proof}

\begin{lemma}[21.8.2]
\label{0.21.8.2-es}
Sean $K$ un cuerpo, $p$ su exponente de característica, $k_0$ un subcuerpo de
$K$, $(k_\alpha)_{\alpha\in I}$ una familia filtrante decreciente de
subcuerpos de $K$ tal que $\bigcap_{\alpha\in I}k_\alpha(K^p)=k_0(K^p)$.
Si $(x_i)_{1\leqslant i\leqslant n}$ es una familia finita de elementos de $K$
que es $p$-independiente sobre $k_0$, existe un índice $\alpha$ tal que $(x_i)$ sea
$p$-independiente sobre $k_\alpha$.

En efecto, decir que la familia $(x_i)$ es $p$-independiente sobre un subcuerpo $k$
de $K$ significa que la familia finita de monomios
$\prod_{i=1}^n x_i^{m(i)}$ con $0\leqslant m(i)<p$ es libre sobre $k(K^p)$;
basta, pues, aplicar el lema (21.8.1) a esta familia de monomios en el
espacio vectorial $V=K$, y a los subcuerpos $k_\alpha(K^p)$ y $k_0(K^p)$
de $K$.
\end{lemma}

\begin{theorem}[21.8.3]
\label{0.21.8.3-es}
Sean $K$ un cuerpo de característica $p>0$, $k_0$ un subcuerpo de $K$,
$(k_\alpha)_{\alpha\in I}$ una familia filtrante decreciente de subcuerpos
de $K$ que contienen a $k_0$. Las condiciones siguientes son equivalentes:
\begin{enumerate}
  \item[\rm a)] $\bigcap_{\alpha\in I}k_\alpha(K^p)=k_0(K^p)$.
  \item[\rm b)] Para toda extensión $L$ de $K$ tal que
  $\Upsilon_{L/K/k_\alpha}$ sea un espacio vectorial sobre $L$ de rango finito
  (lo que ocurre en particular si $L$ es una extensión de tipo finito en
  virtud de (21.7.2)), existe $\alpha\in I$ tal que $k_\alpha$ sea
  $k_0$-admisible para la extensión $L$ de $K$.
  \item[\rm b$'$)] Para toda extensión $L=K(x)$ de $K$, con $x^p\in K$,
  existe $\alpha\in I$ tal que $k_\alpha$ sea $k_0$-admisible para
  la extensión $L$ de $K$.
  \item[\rm c)] La aplicación canónica
\begin{equation}
\tag{21.8.3.1}
\label{0.21.8.3.1-es}
\Omega^1_{K/k_0}\longrightarrow\varprojlim_{\alpha}\Omega^1_{K/k_\alpha}
\end{equation}
es inyectiva.
\end{enumerate}
\end{theorem}

\begin{proof}
La aplicación canónica (21.8.3.1) se obtiene, naturalmente, pasando al límite
proyectivo en el sistema proyectivo de homomorfismos
$\Omega^1_{K/k_0}\to\Omega^1_{K/k_\alpha}$ (20.5.3.3). Probaremos el
teorema según el esquema lógico $c)\Rightarrow b)\Rightarrow b')
\Rightarrow a)\Rightarrow c)$.

Decir que $k_\alpha$ es $k_0$-admisible para la extensión $L$ de $K$ significa
que el homomorfismo
\oldpage[0\textsubscript{IV-1}]{173}
canónico $\Upsilon_{L/K/k_0}\to\Upsilon_{L/K/k_\alpha}$ es inyectivo; ahora
bien, si $N_\alpha$ es el núcleo de este homomorfismo, los $N_\alpha$ forman un
sistema filtrante decreciente de subespacios vectoriales de
$\Upsilon_{L/K/k_0}$, y como $\Upsilon_{L/K/k_0}$ es por hipótesis de rango
finito, equivale a decir que uno de los $N_\alpha$ es nulo o que su
intersección es nula. Pero esta intersección no es sino el núcleo del
homomorfismo límite proyectivo
$\Upsilon_{L/K/k_0}\to\varprojlim_\alpha\Upsilon_{L/K/k_\alpha}$. Ahora bien,
se tiene el diagrama conmutativo
\[
\xymatrix{
\Upsilon_{L/K/k_0}\ar[r]\ar[d]&\varprojlim\Upsilon_{L/K/k_\alpha}\ar[d]\\
\Omega^1_{K/k_0}\otimes_KL\ar[r]&
\varprojlim(\Omega^1_{K/k_\alpha}\otimes_KL)
}
\]
donde la flecha vertical de la izquierda es inyectiva por definición. Para
demostrar que $c)$ implica $b)$, basta, por tanto, mostrar que la aplicación
canónica
\begin{equation}
\tag{21.8.3.2}
\label{0.21.8.3.2-es}
\Omega^1_{K/k_0}\otimes_KV\longrightarrow
\varprojlim_{\alpha}(\Omega^1_{K/k_\alpha}\otimes_KV)
\end{equation}
es inyectiva para todo espacio vectorial $V$ sobre $K$ (y en particular para
$V=L$). Ahora bien, esto es evidente si $V=K^n$, pues entonces
$W\otimes_KV=W^n$ para todo espacio vectorial $W$ sobre $K$ y los productos y
los límites proyectivos conmutan. Por otra parte, para todo elemento $z$ del
primer miembro de (21.8.3.2), existe un subespacio $V'$ de $V$ de rango finito
tal que $z\in\Omega^1_{K/k_0}\otimes_KV'$ (identificado canónicamente con un
subespacio de $\Omega^1_{K/k_0}\otimes_KV$); si $z\ne0$, su imagen en
$\varprojlim_\alpha(\Omega^1_{K/k_\alpha}\otimes_KV')$ no es, por tanto, nula;
como el funtor $\varprojlim$ es exacto por la izquierda, la imagen de $z$ en el
segundo miembro de (21.8.3.2) también es $\ne0$, lo que termina de demostrar
que $c)$ implica $b)$.

Es trivial que $b)$ implica $b')$; mostremos que $b')$ implica $a)$. Con las
notaciones de $b')$, podemos suponer que $L\ne K$. Resulta entonces de (21.4.7)
que se tiene $\operatorname{rg}\Upsilon_{L/K}\leqslant1$. Si
$\Upsilon_{L/K/k_\alpha}=0$, no hay nada que demostrar en virtud de
(21.6.1.1). En caso contrario, $\Upsilon_{L/K/k_\alpha}$ se identifica
canónicamente con $\Upsilon_{L/K}$ y, si se pone $a=x^p$, $\Upsilon_{L/K}$
tiene una base formada por el único elemento $d_K(a)\otimes1$ (21.4.7);
$\Upsilon_{L/K/k_\alpha}$ tiene, por tanto, una base formada por el único
elemento $d_{K/k_\alpha}(a)\otimes1$, y decir que un subcuerpo $k_\alpha$ es
$k_0$-admisible para la extensión $L$ de $K$ significa que se tiene
$d_{K/k_\alpha}(a)\ne0$, o también (21.4.5) que
$a\notin k_\alpha(K^p)$. Ahora bien, para todo $a\notin k_0(K^p)$, se tiene
$d_{K/k_0}(a)\ne0$ y $K(a^{1/p})\ne K$, por lo que se puede aplicar $b')$, que
demuestra la existencia de un $\alpha$ tal que $a\notin k_\alpha(K^p)$; en
otras palabras, $b')$ implica $a)$.

Queda por mostrar que $a)$ implica $c)$. Sea $(x_\mu)_{\mu\in M}$ una
$p$-base de $K$ sobre $k_0$; entonces, los $d_{K/k_0}(x_\mu)$ forman una base
del $K$-espacio vectorial $\Omega^1_{K/k_0}$ (21.4.5), y la condición $c)$
significa que para toda parte finita $J$ del conjunto de índices $M$, existe
un $\alpha\in I$ tal que los $d_{K/k_\alpha}(x_\mu)$, para $\mu\in J$, sean
linealmente independientes en $\Omega^1_{K/k_\alpha}$; pero esto significa
también (21.4.5) que los $x_\mu$, para $\mu\in J$, forman una familia
$p$-libre sobre $k_\alpha$, y la existencia de un $\alpha$ con esta propiedad
resulta de (21.8.2) y de la hipótesis $a)$.
\end{proof}
\oldpage[0\textsubscript{IV-1}]{174}

\begin{corollary}[21.8.4]
\label{0.21.8.4-es}
Sean $K$ un cuerpo de característica $p>0$, $k_0$ un subcuerpo de $K$ y
$(k_\alpha)_{\alpha\in I}$ una familia filtrante decreciente de subcuerpos de
$K$, que contienen a $k_0$, tal que
\[
  \bigcap_{\alpha\in I}k_\alpha(K^p)=k_0(K^p).
\]
Sea $L$ una extensión de $K$ tal que $\Upsilon_{L/K/k_0}$ sea un
$L$-espacio vectorial de rango finito; entonces, para todo cuerpo $K'$ tal que
$K\subset K'\subset L$, existe $\alpha\in I$ tal que $k_\alpha$ sea
$k_0$-admisible para la extensión $L$ de $K'$.
\end{corollary}

\begin{proof}
Basta aplicar (21.8.3) y (21.6.6).
\end{proof}

\begin{corollary}[21.8.5]
\label{0.21.8.5-es}
Sean $K$ un cuerpo de característica $p>0$, $k_0$ un subcuerpo de $K$ y
$(k_\alpha)_{\alpha\in I}$ una familia filtrante decreciente de subcuerpos de
$K$ que contienen a $k_0$, tal que
\[
  \bigcap_{\alpha\in I}k_\alpha(K^p)=k_0(K^p).
\]
Entonces, para toda extensión $L$ de $K$ tal que $\Upsilon_{L/K/k_0}$ sea un
$L$-espacio vectorial de rango finito, se tiene
\[
  \bigcap_{\alpha\in I}k_\alpha(L^p)=k_0(L^p).
\]
\end{corollary}

\begin{proof}
Supongamos, en efecto, que $M$ sea una extensión de $L$ tal que
$\Upsilon_{M/L/k_0}$ sea un $M$-espacio vectorial de rango finito. La sucesión
exacta (21.6.1.1)
\[
  0\longrightarrow \Upsilon_{L/K/k_0}\otimes_L M
  \longrightarrow \Upsilon_{M/K/k_0}
  \longrightarrow \Upsilon_{M/L/k_0}
\]
y la hipótesis muestran entonces que $\Upsilon_{M/K/k_0}$ también es un
$M$-espacio vectorial de rango finito. Existe, por tanto, un índice $\alpha$
tal que $k_\alpha$ sea $k_0$-admisible para la extensión $M$ de $K$ (21.8.3),
y, por consiguiente, también para la extensión $M$ de $L$ (21.6.6); como esto
tiene lugar para toda extensión $M$ de $L$ tal que $\Upsilon_{M/L/k_0}$ sea de
rango finito, el corolario resulta de la equivalencia de $a)$ y $b)$ en
(21.8.3).
\end{proof}

\begin{corollary}[21.8.6]
\label{0.21.8.6-es}
Sean $K$ un cuerpo, $p$ su exponente característico y $k_0$ un subcuerpo de
$K$. Si $L$ es una extensión de $K$ tal que $\Upsilon_{L/K/k_0}$ sea un
$L$-espacio vectorial de rango finito, existe un subcuerpo $k$ de $K$, que
contiene a $k_0(K^p)$, tal que $[K:k]$ sea finito y que sea $k_0$-admisible
para la extensión $L$ de $K$.
\end{corollary}

\begin{proof}
Basta, en efecto, en virtud de (21.8.3), construir una familia filtrante
decreciente $(k_\alpha)_{\alpha\in I}$ de subcuerpos de $K$, que contengan a
$K^p$ y a $k_0$, para los cuales $[K:k_\alpha]<+\infty$ y
$\bigcap_\alpha k_\alpha=k_0(K^p)$. Para ello se considera una $p$-base
$(x_\lambda)_{\lambda\in J}$ de $K$ sobre $k_0$ y, para toda parte finita
$H$ de $J$, se considera el subcuerpo $k_H$ de $K$ engendrado por
$k_0(K^p)$ y los $x_\lambda$ de índice $\lambda\in J-H$; resulta de esta
definición que $(x_\lambda)_{\lambda\in H}$ es una $p$-base de $K$ sobre
$k_H$, y se concluye inmediatamente que los $k_H$ verifican las condiciones
requeridas.
\end{proof}

\begin{remarks}[21.8.7]
\label{0.21.8.7-es}
\begin{enumerate}
  \item[(i)] Ya se ha visto (21.7.2) que si $L$ es una extensión de
  \emph{tipo finito} de $K$, $\Upsilon_{L/K/k_0}$ es de rango finito para
  todo subcuerpo $k_0$ de $K$. Lo mismo sucede si $L$ es una extensión
  \emph{separable} de $K$, pues, en virtud de (20.6.19), se tiene
  $\Upsilon_{L/K/k_0}=0$. Finalmente, si $L$ es una \emph{extensión de tipo
  finito de una extensión separable $L_0$ de $K$}, el mismo razonamiento que
  en (21.8.5) muestra que $\Upsilon_{L/K/k_0}$ sigue siendo de rango finito
  (y de hecho es isomorfo a un subespacio de $\Upsilon_{L/L_0/k_0}$).

  \oldpage[0\textsubscript{IV-1}]{175}
  \item[(ii)] En el enunciado de (21.8.5), y por tanto también en el de
  (21.8.3, $b)$), no se puede omitir la hipótesis de que
  $\Upsilon_{L/K/k_0}$ es de rango finito sobre $L$. Tomemos como $k_0$ el
  cuerpo primo $\mathbf F_p$, y como $K$ un cuerpo tal que $[K:K^p]$ sea
  infinito numerable (por ejemplo, el cuerpo de fracciones racionales
  $\mathbf F_p(X_1,\ldots,X_n,\ldots)$ en una infinidad de indeterminadas); el
  procedimiento de construcción de (21.8.6) muestra inmediatamente que existe
  una sucesión infinita estrictamente decreciente $(K_n)$ de subcuerpos de $K$
  tal que $K_0=K$ y $\bigcap_nK_n=K^p$. Vamos a construir una sucesión
  creciente de extensiones finitas $M_n$ de $K$, tal que, si $M$ es la unión
  de los $M_n$, la extensión $L=M^{1/p}$ de $K$ contradiga (21.8.5). Para
  ello, sea $x$ un elemento de $K-K^p$, pongamos $M_0=K^p$, y
  $M_n=K^p(a_1x,a_2x,\ldots,a_nx)$ para $n\geqslant1$, donde los $a_n$ se
  construyen por recurrencia de modo que $a_n\in K_n$ y
  $a_{n+1}\notin M_n(x)$ para todo $n$: esto es posible, pues $M_n(x)$ es de
  grado finito sobre $K^p$, mientras que $K_n$ no lo es. Se concluye
  inmediatamente que $x\notin M=L^p$, pero como $x=(a_nx)a_n^{-1}$, se tiene
  $x\in K_n(M)=K_n(L^p)$ para todo $n$.
\end{enumerate}
\end{remarks}

\begin{proposition}[21.8.8]
\label{0.21.8.8-es}
Sea $k$ un cuerpo de característica $p>0$, y sean
$A=k[[T_1,\ldots,T_r]]$ el anillo de series formales en $r$ indeterminadas
sobre $k$ y $K=k((T_1,\ldots,T_r))$ su cuerpo de fracciones. Entonces existe
una familia filtrante decreciente $(A_\alpha)$ de subanillos noetherianos de
$A$ tal que $A$ sea un $A_\alpha$-módulo libre de tipo finito para todo
$\alpha$ y que, si $K_\alpha$ es el cuerpo de fracciones de $A_\alpha$, se
tenga $\bigcap_\alpha K_\alpha=K^p$.
\end{proposition}

\begin{proof}
Se puede escribir $K^p=k^p((T_1^p,\ldots,T_r^p))$; se ha visto en la
demostración de (21.8.6) que existe una familia decreciente $(k_\alpha)$ de
subcuerpos de $k$ tal que $[k:k_\alpha]$ sea finito para todo $\alpha$ y
$\bigcap_\alpha k_\alpha=k^p$; está claro que, si se pone
$A_\alpha=k_\alpha[[T_1^p,\ldots,T_r^p]]$, $A$ es un $A_\alpha$-módulo libre
de tipo finito; todo se reduce, por tanto, a demostrar la relación
\begin{equation}
\tag{21.8.8.1}
\label{0.21.8.8.1-es}
  \bigcap_\alpha k_\alpha((T_1^p,\ldots,T_r^p))
  =k^p((T_1^p,\ldots,T_r^p)).
\end{equation}
Como $\bigcap_\alpha k_\alpha=k^p$, esto resultará de los dos lemas siguientes.
\end{proof}

\begin{lemma}[21.8.8.2]
\label{0.21.8.8.2-es}
Si $k'$ es una extensión de un cuerpo $k$, se tiene
\[
  k'[[T_1,\ldots,T_r]]\cap k((T_1,\ldots,T_r))=k[[T_1,\ldots,T_r]].
\]
\end{lemma}

\begin{proof}
En efecto, pongamos $C=k[[T_1,\ldots,T_r]]$ y
$D=k'[[T_1,\ldots,T_r]]$; como $k((T_1,\ldots,T_r))$ es el cuerpo de
fracciones de $C$, bastará demostrar que $D$ es un $C$-módulo
\emph{fielmente plano} (Bourbaki, \emph{Alg. comm.}, cap.~I, \S~3,
n\textsuperscript{o}~5, prop.~10). Ahora bien, $C$ y $D$ son anillos locales
noetherianos, y si $\mathfrak m$ es el ideal maximal de $C$, se tiene
$D/\mathfrak m^jD=(C/\mathfrak m^j)\otimes_k k'$, por lo que
$D/\mathfrak m^jD$ es un $(C/\mathfrak m^j)$-módulo plano; basta, por tanto,
aplicar ($\mathbf 0_{\mathrm{III}}$, 10.2.1) y
($\mathbf 0_{\mathrm I}$, 6.6.2).
\end{proof}

\begin{lemma}[21.8.8.3]
\label{0.21.8.8.3-es}
Sean $k$ un cuerpo, $(k_\alpha)$ una familia filtrante decreciente de
subcuerpos de $k$, y pongamos $k_0=\bigcap_\alpha k_\alpha$. Se supone que
existe una potencia $q$ del exponente característico de $k$ tal que
$k^q\subset k_0$. Entonces se tiene
\begin{equation}
\tag{21.8.8.4}
\label{0.21.8.8.4-es}
  \bigcap_\alpha k_\alpha((T_1,\ldots,T_r))=k_0((T_1,\ldots,T_r)).
\end{equation}
\end{lemma}

\begin{proof}
\oldpage[0\textsubscript{IV-1}]{176}
Basta demostrar que un elemento $f\ne0$ del primer miembro de (21.8.8.4)
pertenece al segundo miembro. Sean $\gamma$ un índice y
$g\in k_\gamma[[T_1,\ldots,T_r]]$ tales que
$gf\in k_\gamma[[T_1,\ldots,T_r]]$; reemplazando, si es necesario, $g$ por
$g^q$, se puede suponer que $g\in k_0[[T_1,\ldots,T_r]]$. Entonces, para todo
$\alpha\geqslant\gamma$, se tiene
\[
  gf\in k_\gamma[[T_1,\ldots,T_r]]\cap k_\alpha((T_1,\ldots,T_r))
  =k_\alpha[[T_1,\ldots,T_r]]
\]
en virtud del lema (21.8.8.2). Pero está claro que la intersección de los
anillos $k_\alpha[[T_1,\ldots,T_r]]$ no es sino
$k_0[[T_1,\ldots,T_r]]$, y se tiene, por tanto,
$f\in k_0((T_1,\ldots,T_r))$.
\end{proof}

\subsection{Módulos de diferenciales completados en los anillos de series formales}
\label{subsection:0.21.9-es}

En este número, ya no se supone necesariamente que los cuerpos sean de
característica $>0$.

\begin{lemma}[21.9.1]
\label{0.21.9.1-es}
Sean $k$ un cuerpo, $A$ un anillo local noetheriano completo que es una
$k$-álgebra y $K$ el cuerpo residual de $A$. Si $K$ es una extensión de tipo
finito de $k$, el $A$-módulo $\widehat\Omega^1_{A/k}$ es de tipo finito.
\end{lemma}

\begin{proof}
En virtud de (20.7.15), basta demostrar que $\Omega^1_{K/k}$ es un
$K$-espacio vectorial de rango finito. Ahora bien, por hipótesis, $K$ es el
cuerpo de fracciones de una $k$-álgebra de tipo finito $B$; como
$\Omega^1_{B/k}$ es un $B$-módulo de tipo finito en virtud de (20.4.7), la
conclusión resulta de (20.5.9).
\end{proof}

\begin{proposition}[21.9.2]
\label{0.21.9.2-es}
Sean $k_0$ un cuerpo, $A$ un anillo local noetheriano completo que es una
$k_0$-álgebra formalmente lisa, $\mathfrak m$ su ideal maximal y
$K=A/\mathfrak m$ su cuerpo residual (de modo que, como anillo, $A$ es
isomorfo a un anillo de series formales $K[[T_1,\ldots,T_r]]$ (19.6.5)).
Si $K$ es una extensión de tipo finito de $k_0$, entonces
$\widehat\Omega^1_{A/k_0}$ es un $A$-módulo libre de rango igual a
$\dim(A)+\deg.\operatorname{tr}_{k_0}K$. Además, para todo subcuerpo $k$ de
$k_0$ tal que $\Omega^1_{k_0/k}$ sea de rango finito,
$\widehat\Omega^1_{A/k}$ es un $A$-módulo libre de rango igual a
\[
  \dim(A)+\deg.\operatorname{tr}_{k_0}K
  +\operatorname{rg}_{k_0}(\Omega^1_{k_0/k}).
\]
\end{proposition}

\begin{proof}
Está claro que, si $P$ es el subcuerpo primo de $k_0$, $K$ es una
$P$-álgebra formalmente lisa (19.6.1). Observemos, por otra parte, que se tiene
$\Upsilon^K_{A/k_0/P}=0$: en efecto, esto equivale a ver que el homomorfismo
\[
  u_{A/k_0/P}\otimes1_K:
  \Omega^1_{k_0/P}\otimes_{k_0}K
  \longrightarrow \Omega^1_{A/P}\otimes_A K
\]
es inyectivo (20.6.14.5), donde $u$ designa el homomorfismo estructural
$k_0\to A$. Pero como $A$ es una $k_0$-álgebra formalmente lisa por hipótesis,
resulta de (20.7.2) que $u_{A/k_0/P}$ es formalmente invertible por la
izquierda, por lo que (19.1.7) $u_{A/k_0/P}\otimes1_K$ es invertible por la
izquierda y, \emph{a fortiori}, inyectivo, lo que demuestra nuestra afirmación.
Aplicando la sucesión exacta (20.6.22.1), donde $\Lambda,A,B,C$ se sustituyen
por $P,k_0,A,K$, se obtiene la sucesión exacta
\[
  0\longrightarrow \Upsilon^K_{A/k_0}
  \longrightarrow \mathfrak m/\mathfrak m^2
  \longrightarrow \Omega^1_{A/k_0}\otimes_A K
  \longrightarrow \Omega^1_{K/k_0}\longrightarrow0,
\]
de modo que se tiene
\[
\begin{split}
  \operatorname{rg}_K(\Omega^1_{A/k_0}\otimes_AK)
  &={\operatorname{rg}}_K(\mathfrak m/\mathfrak m^2)
    +\bigl(\operatorname{rg}_K(\Omega^1_{K/k_0})
    -\operatorname{rg}_K(\Upsilon_{K/k_0})\bigr)\\
  &={\operatorname{rg}}_K(\mathfrak m/\mathfrak m^2)
    +\deg.\operatorname{tr}_{k_0}K
\end{split}
\]
en virtud de (21.7.1). Por otra parte, $\widehat\Omega^1_{A/k_0}$ es un
$A$-módulo formalmente proyectivo (20.4.9), y como su topología es la topología
$\mathfrak m$-preádica (20.4.5),
$\widehat\Omega^1_{A/k_0}/\mathfrak m^{j+1}\widehat\Omega^1_{A/k_0}$ es un
$(A/\mathfrak m^{j+1})$-módulo proyectivo para todo $j$ (19.2.4), luego libre
($\mathbf 0_{\mathrm{III}}$, 10.1.2) y de rango

\oldpage[0\textsubscript{IV-1}]{177}
$n=\operatorname{rg}_K(\Omega^1_{A/k_0}\otimes_AK)$. En virtud de (21.9.1),
$\widehat\Omega^1_{A/k_0}$ es un $A$-módulo de tipo finito y, además, este
$A$-módulo es plano en virtud de ($\mathbf 0_{\mathrm{III}}$, 10.2.1); es,
por tanto, libre ($\mathbf 0_{\mathrm{III}}$, 10.1.3) de rango $n$, de donde
resulta la primera afirmación.

Como por hipótesis $\Omega^1_{k_0/k}$ es de rango finito sobre $k_0$, el
producto tensorial completado
$\Omega^1_{k_0/k}\widehat\otimes_{k_0}A$ se identifica con
$\Omega^1_{k_0/k}\otimes_{k_0}A$; como $A$ es una $k_0$-álgebra formalmente
lisa, resulta de (20.7.18) que el homomorfismo
\[
  \widehat u_{A/k_0/k}:
  \Omega^1_{k_0/k}\otimes_{k_0}A
  \longrightarrow\widehat\Omega^1_{A/k}
\]
admite un inverso por la izquierda que es un $A$-homomorfismo continuo, lo que
implica en particular que su imagen es cerrada en
$\widehat\Omega^1_{A/k}$; la sucesión exacta
\[
  0\longrightarrow \Omega^1_{k_0/k}\otimes_{k_0}A
  \xrightarrow{\widehat u_{A/k_0/k}}\widehat\Omega^1_{A/k}
  \longrightarrow\widehat\Omega^1_{A/k_0}\longrightarrow0
\]
es, por tanto, exacta y escindida en virtud de (20.7.17) y de lo que precede;
esto termina de demostrar la proposición.
\end{proof}

\begin{corollary}[21.9.3]
\label{0.21.9.3-es}
Sean $k$ un cuerpo y $A$ el anillo de series formales
$k[[T_1,\ldots,T_r]]$, provisto de su estructura usual de $k$-álgebra.
Entonces $\widehat\Omega^1_{A/k}$ es un $A$-módulo libre de rango
$r=\dim(A)$.
\end{corollary}

\begin{proof}
Basta observar que $A$ es una $k$-álgebra formalmente lisa (19.3.4), y aplicar
(21.9.2) con $k_0=k=K$.
\end{proof}

\begin{lemma}[21.9.4]
\label{0.21.9.4-es}
Sean $k$ un cuerpo de característica $p>0$, $A$ una $k$-álgebra que es un
anillo local noetheriano completo y $B$ una sub-$k$-álgebra de $A$ isomorfa a
una $k$-álgebra topológica de la forma $k[[T_1,\ldots,T_r]]$ y tal que $A$ sea
una $B$-álgebra de tipo finito. Si $B_1$ es la subálgebra
$k[[T_1^p,\ldots,T_r^p]]$ de $B$, $\widehat\Omega^1_{A/k}$ se identifica
canónicamente con $\Omega^1_{A/B_1}$.
\end{lemma}

\begin{proof}
Toda derivación continua de $B_1$ en un $A$-módulo topológico $P$ que sea la
restricción de una $k$-derivación continua de $A$ en $P$ es \emph{nula}, pues
lo es en el subanillo $k[T_1^p,\ldots,T_r^p]$ de los polinomios, y este último
es denso en $B_1$. La sucesión exacta (20.3.7.1) muestra, por tanto, que el
homomorfismo canónico
\[
  \operatorname{D\acute{e}r.\,cont}_{B_1}(A,P)
  \longrightarrow \operatorname{D\acute{e}r.\,cont}_k(A,P)
\]
es biyectivo; la conclusión resulta de (20.7.14.4), teniendo en cuenta que
$\Omega^1_{A/B_1}$ es un $A$-módulo de tipo finito (20.4.7), luego separado y
completo, puesto que $A$ es un anillo local noetheriano completo.
\end{proof}

\begin{proposition}[21.9.5]
\label{0.21.9.5-es}
Sean $A$ un anillo local noetheriano completo íntegro y $A_0$ un subanillo de
$A$ isomorfo a un anillo de series formales $k[[T_1,\ldots,T_r]]$ sobre un
cuerpo $k$, tal que $A$ sea una $A_0$-álgebra finita y que el cuerpo de
fracciones $E$ de $A$ sea una extensión separable del cuerpo de fracciones
$L_0$ de $A_0$. Entonces se tiene
\begin{equation}
\tag{21.9.5.1}
\label{0.21.9.5.1-es}
  \operatorname{rg}_A\widehat\Omega^1_{A/k}
  =\operatorname{rg}_E(\widehat\Omega^1_{A/k}\otimes_AE)=\dim(A).
\end{equation}
\end{proposition}

\begin{proof}
Observemos que si $\mathfrak m_0$ es el ideal maximal de $A_0$, la topología
de $A$ es la topología $\mathfrak m_0$-ádica, puesto que $A$ es una
$A_0$-álgebra finita (Bourbaki, \emph{Alg. comm.}, cap.~IV, \S~2,
n\textsuperscript{o}~5, cor.~3 de la prop.~9), e induce en $A_0$ la topología
$\mathfrak m_0$-ádica (Bourbaki, \emph{Alg. comm.}, cap.~III, \S~3,
n\textsuperscript{o}~4, teo.~3). Se sabe (21.9.1) que
$\widehat\Omega^1_{A/k}$ es un $A$-módulo de tipo finito y,

\oldpage[0\textsubscript{IV-1}]{178}
en virtud de (21.9.3), $\widehat\Omega^1_{A_0/k}$ es un $A_0$-módulo libre de
rango $r=\dim(A_0)=\dim(A)$ (16.1.5); el producto tensorial
$\widehat\Omega^1_{A_0/k}\otimes_{A_0}A$ es, por tanto, un $A$-módulo libre de
rango $r$ idéntico a su completado separado; finalmente,
$\widehat\Omega^1_{A/A_0}$ es un $A$-módulo de tipo finito (20.4.7), luego
idéntico a su completado separado ($\mathbf 0_{\mathrm I}$, 7.3.6). Sentado
esto, en la sucesión de homomorfismos
\begin{equation}
\tag{21.9.5.2}
\label{0.21.9.5.2-es}
  \widehat\Omega^1_{A_0/k}\otimes_{A_0}A
  \xrightarrow{u}\widehat\Omega^1_{A/k}
  \xrightarrow{v}\widehat\Omega^1_{A/A_0}\longrightarrow0
\end{equation}
se sabe que $v$ es sobreyectivo y que $\operatorname{Im}(u)$ es densa en
$\operatorname{Ker}(v)$ (20.7.17); pero todo sub-$A_0$-módulo de
$\widehat\Omega^1_{A/k}$ es cerrado para la topología $\mathfrak m_0$-ádica
($\mathbf 0_{\mathrm I}$, 7.3.5), por lo que la sucesión (21.9.5.2) es exacta.
Observemos, por otra parte, que como $A$ es entero sobre $A_0$, $A_0$ es
íntegramente cerrado y $E$ es separable sobre $L_0$,
$\widehat\Omega^1_{A/A_0}$ es un $A$-módulo de torsión (20.4.13, (iv)); al
tensorizar la sucesión exacta (21.9.5.2) por $E$, se obtiene una sucesión
exacta
\[
  \widehat\Omega^1_{A_0/k}\otimes_{A_0}E
  \longrightarrow \widehat\Omega^1_{A/k}\otimes_AE\longrightarrow0
\]
de donde $\operatorname{rg}_E(\widehat\Omega^1_{A/k}\otimes_AE)\leqslant r$.
Queda, por tanto, mostrar que existen, en
$\widehat\Omega^1_{A/k}\otimes_AE$, $r$ elementos linealmente independientes
sobre $E$. Ahora bien, sean $D_i=\partial/\partial T_i$ las derivaciones
canónicas de $A_0$ en sí mismo $(1\leqslant i\leqslant r)$; se prolongan de
manera única a derivaciones (también denotadas $D_i$) de $E$ en sí mismo,
puesto que $E$ es una extensión separable finita de $L_0$; por restricción a
$A$, estas derivaciones dan derivaciones de $A$ en $E$, y es inmediato que
estas derivaciones toman sus valores en un mismo sub-$A$-módulo de tipo finito
de $E$; como son continuas en $A_0$ y la topología de $A$ es la topología
$\mathfrak m_0$-ádica, se han formado así $r$ derivaciones continuas de $A$
en un $A$-módulo topológico, que son evidentemente linealmente independientes,
puesto que $D_i(T_j)=\delta_{ij}$; esto termina de demostrar la proposición.
\end{proof}

\begin{proposition}[21.9.6]
\label{0.21.9.6-es}
Sean $k_0\subset k\subset K_0$ tres cuerpos de característica $p>0$, tales
que $[K_0:k]<+\infty$; se pone
\[
  L_0=K_0((T_1,\ldots,T_r)),\qquad
  M_0=K_0((T_1^p,\ldots,T_r^p)),
\]
\[
  L=k((T_1,\ldots,T_r)),\qquad
  M=k((T_1^p,\ldots,T_r^p)),\qquad
  N=k((T_1^{p^2},\ldots,T_r^{p^2})).
\]
Sea $E$ una extensión de tipo finito de $L_0$.
\begin{enumerate}
  \item[(i)] Si $\Upsilon_{K_0/k_0}$ es de rango finito, se tiene
  \begin{equation}
  \tag{21.9.6.1}
  \label{0.21.9.6.1-es}
    \operatorname{rg}_E\Omega^1_{E/M}
    -\operatorname{rg}_{K_0}\Omega^1_{K_0/k}
    =(r+\deg.\operatorname{tr}_{L_0}E)
      +d(E/K_0,k_0)+d(E/M_0,N/k_0)
  \end{equation}
  (notaciones de (21.6.1)), y en particular se tiene
  \begin{equation}
  \tag{21.9.6.2}
  \label{0.21.9.6.2-es}
    \operatorname{rg}_E\Omega^1_{E/M}
    -\operatorname{rg}_{K_0}\Omega^1_{K_0/k}
    \geqslant r+\deg.\operatorname{tr}_{L_0}E+d(E/K_0,k_0).
  \end{equation}
  Además, si los dos miembros de (21.9.6.2) son iguales, siguen siéndolo al
  sustituir $k$ por un cuerpo $k'$ tal que $k_0\subset k'\subset k$.

  \item[(ii)] Se supone además que $[k_0:k_0^p]<+\infty$. Sea
  $(k_\alpha)$ una familia filtrante decreciente de subcuerpos de $K_0$ que
  contienen a $k_0$, tales que $[K_0:k_\alpha]<+\infty$ para todo $\alpha$ y
  que $\bigcap_\alpha k_\alpha(K_0^p)=k_0(K_0^p)$; entonces existe un
  $\alpha$ tal que, para $k=k_\alpha$, los dos miembros de (21.9.6.2) sean
  iguales.
\end{enumerate}
\end{proposition}

\begin{proof}
\begin{enumerate}
  \oldpage[0\textsubscript{IV-1}]{179}
  \item[(i)] Se sabe (21.1.5, (ii)) que
  $\operatorname{rg}_E\Omega^1_{E/M}$ no cambia cuando se sustituye $M$ por
  $M(L_0^p)$; como
  $L_0^p=K_0^p((T_1^p,\ldots,T_r^p))$ y $[k(K_0^p):k]<+\infty$, se tiene
  \[
    M(L_0^p)=k(K_0^p)((T_1^p,\ldots,T_r^p))~;
  \]
  asimismo, $\Omega^1_{K_0/k}$ no cambia cuando en él se sustituye $k$ por
  $k(K_0^p)$, de modo que se puede \emph{suponer que $K_0^p\subset k$}, como
  haremos en todo lo que sigue. Introduciremos también los cuerpos
  \[
    k'_0=k_0(K_0^p),\qquad
    N'=k'_0((T_1^{p^2},\ldots,T_r^{p^2}))
  \]
  de modo que se tiene el diagrama de cuerpos
  \[
    \xymatrix{
      K_0\ar[r] & M_0\ar[r] & L_0\ar[r] & E\\
      k\ar[u]\ar[r] & N\ar[u]\ar[r] & M\ar[u]\ar[r] & L\ar[u]\\
      k_0\ar[u]\ar[r] & k'_0\ar[u]\ar[r] & N'\ar[u] &
    }
  \]

  Nos proponemos evaluar la diferencia $\delta$ entre el primer y el segundo
  miembro de (21.9.6.2).

  Resulta de (21.7.3) que se tiene
  \[
    \operatorname{rg}_E\Omega^1_{E/M}
    -\operatorname{rg}_{M_0}\Omega^1_{M_0/M}
    =\deg.\operatorname{tr}_{M_0}E+d(E/M_0,M)
  \]
  y como $L_0$ es una extensión finita de $M_0$, se tiene
  $\deg.\operatorname{tr}_{L_0}E=\deg.\operatorname{tr}_{M_0}E$. Por otra
  parte, una $p$-base de $K_0$ sobre $k=k(K_0^p)$ es también una $p$-base de
  $M_0$ sobre $M$, por lo que (21.4.5) se tiene
  \[
    \operatorname{rg}_{M_0}\Omega^1_{M_0/M}
    =\operatorname{rg}_{K_0}\Omega^1_{K_0/k}
  \]
  y, por consiguiente,
  \begin{equation}
  \tag{21.9.6.3}
  \label{0.21.9.6.3-es}
    \delta=d(E/M_0,M)-d(E/K_0,k_0)-r.
  \end{equation}

  Observemos ahora el lema clásico:

  \begin{lemma}[21.9.6.4]
  \label{0.21.9.6.4-es}
  Para todo cuerpo $k$, el cuerpo de series formales
  $K=k((T_1,\ldots,T_r))$ es una extensión separable de $k$.
  \end{lemma}

  \begin{proof}
  Recordemos brevemente la demostración de este lema, para ser completos.
  Basta (Bourbaki, \emph{Alg.}, cap.~VIII, \S~7,
  n\textsuperscript{o}~3, demostración del teo.~1) demostrar que, para toda
  extensión finita $k'$ de $k$, $K\otimes_k k'$ no tiene elementos nilpotentes;
  pero si se pone $A=k[[T_1,\ldots,T_r]]$ y $S=A-\{0\}$,
  $K\otimes_k k'$ es igual a $S^{-1}(A\otimes_k k')$, y
  $A\otimes_k k'$ se identifica canónicamente con el anillo íntegro
  $A'=k'[[T_1,\ldots,T_r]]$, y $A$ con un subanillo de $A'$, de donde resulta
  la conclusión.
  \end{proof}

  Utilizando este lema y (21.6.2, (ii)), se tiene $d(L_0/K_0,k_0)=0$, y la
  fórmula (21.6.4.4) da, por tanto,
  \[
    d(E/K_0,k_0)=d(E/L_0,k_0).
  \]

  \oldpage[0\textsubscript{IV-1}]{180}
  Las fórmulas (21.6.4.4) y (21.6.4.2) dan, por otra parte,
  \[
    d(E/M_0,M)=d(E/L_0,M)+d(L_0/M_0,M),
  \]
  \[
    d(L_0/M_0,M)=d(L_0/M_0,M/N)+d(L_0/M_0,N),
  \]
  de donde
  \[
    \delta=d(E/L_0,M)+d(L_0/M_0,N)-d(E/L_0,k_0)
    +d(L_0/M_0,M/N)-r.
  \]

  Vamos a demostrar que se tiene la relación
  \begin{equation}
  \tag{21.9.6.5}
  \label{0.21.9.6.5-es}
    d(L_0/M_0,M/N)=r.
  \end{equation}
  Para ello, observemos que, por definición, se tiene
  \[
    d(L_0/M_0,M/N)
    =\operatorname{rg}_{L_0}\Upsilon_{L_0/M/N}
    -\operatorname{rg}_{M_0}\Upsilon_{M_0/M/N}.
  \]
  Teniendo en cuenta las dos sucesiones exactas
  \[
  \begin{aligned}
    0&\longrightarrow\Upsilon_{L_0/M/N}
      \longrightarrow\Omega^1_{M/N}\otimes_ML_0
      \longrightarrow\Omega^1_{L_0/N}
      \longrightarrow\Omega^1_{L_0/M}\longrightarrow0,\\
    0&\longrightarrow\Upsilon_{M_0/M/N}
      \longrightarrow\Omega^1_{M/N}\otimes_MM_0
      \longrightarrow\Omega^1_{M_0/N}
      \longrightarrow\Omega^1_{M_0/M}\longrightarrow0,
  \end{aligned}
  \]
  se obtiene
  \[
  \begin{split}
    d(L_0/M_0,M/N)={}&\operatorname{rg}_{L_0}\Omega^1_{L_0/M}
      -\operatorname{rg}_{L_0}\Omega^1_{L_0/N}\\
      &-\operatorname{rg}_{M_0}\Omega^1_{M_0/M}
      +\operatorname{rg}_{M_0}\Omega^1_{M_0/N}.
  \end{split}
  \]
  Ahora bien, se tiene $N(L_0^p)\supset M$, luego
  $\Omega^1_{L_0/N}=\Omega^1_{L_0/M}$ (21.1.5, (ii)). Por otra parte, los
  $T_i^p$ $(1\leqslant i\leqslant r)$ forman una $p$-base de $M$ sobre $N$,
  y una $p$-base de $K_0$ sobre $k$ es también una $p$-base de $M_0$ sobre
  $M$, por lo que, como $M_0^p\subset N$,
  \[
    \operatorname{rg}_{M_0}\Omega^1_{M_0/N}
    =r+\operatorname{rg}_{K_0}\Omega^1_{K_0/k}
    =r+\operatorname{rg}_{M_0}\Omega^1_{M_0/M}
  \]
  en virtud de (21.4.5) y (21.1.10), lo que demuestra (21.9.6.5).

  Se tiene además, por (21.6.4.2),
  \[
    d(E/L_0,M)=d(E/L_0,k_0)+d(E/L_0,M/k_0)
  \]
  y
  \[
    d(L_0/M_0,N)=d(L_0/M_0,N/k)+d(L_0/M_0,k).
  \]
  Pero como $L_0$ es separable sobre $K_0$ (21.9.6.4), se tiene
  $d(L_0/K_0,k)=0$, y por consiguiente también $d(L_0/M_0,k)=0$ ((21.6.2) y
  (21.6.4)). Se obtiene, por tanto,
  \[
    \delta=d(E/L_0,M/k_0)+d(L_0/M_0,N/k).
  \]
  Apliquemos de nuevo dos veces (21.6.4.2) al primer término del segundo
  miembro de esta fórmula; se obtiene
  \[
    d(E/L_0,M/k_0)=d(E/L_0,M/N)+d(E/L_0,N/k)+d(E/L_0,k/k_0).
  \]
  Como $M\subset N(L_0^p)$, se tiene $d(E/L_0,M/N)=0$ (21.6.2), luego (por
  (21.6.4.4))
  \[
    \delta=d(E/L_0,k/k_0)+d(E/M_0,N/k).
  \]
  Pero $M_0$ y $L_0$ son extensiones separables de $K_0$ (21.9.6.4), por lo
  que
  \[
    d(L_0/K_0,k/k_0)=d(M_0/K_0,k/k_0)=0
  \]
  (21.6.2), y aplicando de nuevo dos veces (21.6.4.4),

  \oldpage[0\textsubscript{IV-1}]{181}
  se tiene
  \[
    d(E/L_0,k/k_0)=d(E/K_0,k/k_0)=d(E/M_0,k/k_0),
  \]
  de donde, por una última aplicación de (21.6.4.2),
  \[
    \delta=d(E/M_0,N/k_0),
  \]
  lo que demuestra (21.9.6.1). Además, cuando $k$ se sustituye por una
  subextensión $k'$ de $k_0$, $N$ se sustituye por una subextensión $N'$, por
  lo que $d(E/M_0,N'/k_0)\leqslant d(E/M_0,N/k_0)$ (21.6.4.2), lo que
  demuestra la última afirmación de (i).

  \item[(ii)] Para todo $\alpha$, pongamos
  $N_\alpha=k_\alpha((T_1^{p^2},\ldots,T_r^{p^2}))$. En virtud de (21.9.6.1)
  y de (21.8.3), se trata de demostrar (teniendo en cuenta que
  $K_0^p\subset k_\alpha$ y que $E$ es una extensión de tipo finito de $M_0$)
  que se tiene
  \[
    \bigcap_\alpha N_\alpha=k_0(M_0^p).
  \]
  Ahora bien, se tiene
  \begin{equation}
  \tag{21.9.6.6}
  \label{0.21.9.6.6-es}
    \bigcap_\alpha N_\alpha=N'
    =(k_0(K_0^p))((T_1^{p^2},\ldots,T_r^{p^2}))
  \end{equation}
  en virtud de (21.8.8.3) y de la relación
  $\bigcap_\alpha k_\alpha=k_0(K_0^p)$ (21.8.6). Por otra parte, se tiene
  $k_0(M_0^p)=k_0(K_0^p((T_1^{p^2},\ldots,T_r^{p^2})))$. Para terminar la
  demostración de (21.9.6, (ii)), queda, por tanto, demostrar el

  \begin{lemma}[21.9.6.7]
  \label{0.21.9.6.7-es}
  Sean $k_0$ un cuerpo de característica $p>0$ tal que
  $[k_0:k_0^p]<+\infty$ y $K_0$ una extensión de $k_0$. Entonces se tiene
  \begin{equation}
  \tag{21.9.6.8}
  \label{0.21.9.6.8-es}
    (k_0(K_0^p))((T_1,\ldots,T_r))
    =k_0\bigl(K_0^p((T_1,\ldots,T_r))\bigr).
  \end{equation}
  \end{lemma}

  \begin{proof}
  Como $k_0^p\subset K_0^p$, se tiene
  \[
    k_0^p\bigl(K_0^p((T_1,\ldots,T_r))\bigr)
    =K_0^p((T_1,\ldots,T_r)).
  \]
  Si $(c_i)_{1\leqslant i\leqslant n}$ es una base de $k_0$ sobre $k_0^p$,
  es inmediato que $(c_i)$ es también un sistema de generadores de cada uno
  de los dos miembros de (21.9.6.8) sobre
  $K_0^p((T_1,\ldots,T_r))$.
  \end{proof}
\end{enumerate}
\end{proof}

\begin{remark}[21.9.7]
\label{0.21.9.7-es}
La desigualdad (21.9.6.2) puede ser estricta si se omite la hipótesis de que
$[k_0:k_0^p]$ es finito. Tomemos, en efecto, como $k_0$ el cuerpo
$\mathbf F_p(X_n)_{n\geqslant0}$ de las fracciones racionales en una infinidad
de indeterminadas, de modo que los $X_n$ formen una $p$-base de $k_0$ sobre
$\mathbf F_p$, y por consiguiente $[k_0:k_0^p]=+\infty$. En el enunciado de
(21.9.6), tomemos $K_0=k_0$ (por tanto, necesariamente $k=k_0$),
$L_0=k_0((T))$, y como $E$ la extensión $L_0(y)$, donde $y$ es raíz del
polinomio
\[
  Y^p-\sum_{n=0}^{\infty}X_nT^{np}
\]
de $L_0[Y]$. Entonces la diferencia $\delta$ entre los dos miembros de
(21.9.6.2) es \emph{no nula}. En efecto, $E$ es \emph{separable} sobre $k_0$,
pues de lo contrario $y^p$ sería elemento de $k_0(L_0^p)$, y se ve
inmediatamente que esto no es posible (debido a que hay una infinidad de
$X_n$ algebraicamente independientes); se tiene, por tanto,
$d(E/K_0,k_0)=0$, la fórmula (21.9.6.3) se reduce a
$\delta=d(E/M_0,M_0)-1$, y está claro que
$d(E/M_0,M_0)=\operatorname{rg}_E\Upsilon_{E/M_0}$. Todo se reduce a comprobar
que este rango es $2$. Ahora bien, se tiene $E^p\subset M_0$, una $p$-base de
$E$ sobre $L_0$ está formada por el único elemento $y$, y $L_0$ tiene
evidentemente una $p$-base sobre $M_0$ formada por el único elemento $T$; por
tanto, como $E$ es una extensión algebraica de $M_0$, nuestra afirmación
resulta de (21.7.1) y (21.4.5).
\end{remark}

\begin{proposition}[21.9.8]
\label{0.21.9.8-es}
Sean $k_0$ un cuerpo (de característica cualquiera), $K_0$ una extensión
separable de $k_0$ y $A$ un anillo local noetheriano completo, que es una
$k_0$-álgebra y cuyo cuerpo residual

\oldpage[0\textsubscript{IV-1}]{182}
es una extensión finita de $K_0$. Para todo ideal primo $\mathfrak p$ de $A$,
sea $k(\mathfrak p)$ el cuerpo residual de $A$ en $\mathfrak p$. Entonces,
para todo cuerpo $k$ tal que $k_0\subset k\subset K_0$ y tal que
$[K_0:k]<+\infty$, se tiene
\begin{equation}
\tag{21.9.8.1}
\label{0.21.9.8.1-es}
\begin{split}
  &\operatorname{rg}_{k(\mathfrak p)}
    \bigl(\widehat\Omega^1_{(A/\mathfrak p)/k}
      \otimes_{A/\mathfrak p}k(\mathfrak p)\bigr)
    -\operatorname{rg}_{K_0}\Omega^1_{K_0/k}\\
  &\hspace{35mm}\geqslant
    \dim(A/\mathfrak p)
    +\operatorname{rg}_{k(\mathfrak p)}\Upsilon_{k(\mathfrak p)/k_0}.
\end{split}
\end{equation}
Supongamos además que $[k_0:k_0^p]<+\infty$ (donde $p$ es el exponente
característico de $k_0$). Entonces se puede encontrar un cuerpo $k$ tal que
$k_0(K_0^p)\subset k\subset K_0$ y $[K_0:k]<+\infty$, para el cual los dos
miembros de (21.9.8.1) son iguales.
\end{proposition}

\begin{proof}
Como $A/\mathfrak p$ tiene el mismo cuerpo residual $K$ que $A$ y es completo,
se puede reducir al caso en que $A$ es íntegro y $\mathfrak p=0$; se pondrá
entonces $E=k(\mathfrak p)$, cuerpo de fracciones de $A$. Como $K_0$ es
formalmente liso sobre $k_0$ (19.6.1), existe un $k_0$-monomorfismo
$K_0\to A$ que hace de $A$ una $K_0$-álgebra; se sabe además que existe una
sub-$K_0$-álgebra $A_0$ de $A$, $K_0$-isomorfa a un álgebra de series formales
$K_0[[T_1,\ldots,T_r]]$ y tal que $A$ sea una $A_0$-álgebra \emph{finita}
(19.8.9), lo que implica que $E$ es una extensión finita del cuerpo de
fracciones $L_0=K_0((T_1,\ldots,T_r))$ de $A_0$. Se tiene entonces
$\Upsilon_{K_0/k_0}=0$, puesto que $K_0$ es separable sobre $k_0$, y
$\deg.\operatorname{tr}_{L_0}E=0$; además, como $A$ es una $A_0$-álgebra
finita, se tiene $\dim(A)=\dim(A_0)$ (16.1.5).

Si $p=1$, se tiene
$\operatorname{rg}_E(\widehat\Omega^1_{A/k}\otimes_AE)=r$ por (21.9.5) y
(17.1.4, (iii)), y $\Omega^1_{K_0/k}=0$, puesto que $K_0$ es una extensión
finita separable de $k$ (21.7.4); por otra parte,
$\Upsilon_{E/k_0}=0$, ya que $E$ es una extensión separable de $k_0$; los dos
miembros de (21.9.8.1) son, por tanto, iguales en este caso.

Supongamos ahora que $p>1$; si $B=k[[T_1,\ldots,T_r]]$, $A$ es una
$B$-álgebra finita, por lo que, poniendo
$B_1=k[[T_1^p,\ldots,T_r^p]]$, $\widehat\Omega^1_{A/k}$ se identifica con
$\Omega^1_{A/B_1}$ (21.9.4), y, designando por $M$ el cuerpo de fracciones
$k((T_1^p,\ldots,T_r^p))$, resulta de (20.5.9) que
$\widehat\Omega^1_{A/k}\otimes_AE$ se identifica con $\Omega^1_{E/M}$; la
desigualdad (21.9.8.1) no es entonces sino (21.9.6.2), y la última afirmación
de la proposición resulta de (21.9.6, (ii)).\footnote{La fuente francesa
imprime «del corolario»; el resultado que se está demostrando es la
proposición (21.9.8).}
\end{proof}
